% Generated by build.py from paper.template.html and the formulas in build.py.
% Do not edit this file. Edit the template or build.py and run build.py again.
% Compile with XeLaTeX or LuaLaTeX (the file uses fontspec), for example: tectonic falafel-paper.tex
\documentclass[11pt,a4paper]{amsart}
\usepackage[text={5.9in,9.55in},centering]{geometry}
\raggedbottom  % a page that ends early stays short; its gaps are not stretched to fill it
\usepackage{fontspec}
\newfontfamily\leanfont[NFSSFamily=leanmono]{DejaVuSansMono.ttf}
\usepackage{amssymb}
\usepackage{array,booktabs,float}
\usepackage{fancyvrb}
\usepackage{tikz}
\definecolor{accent}{HTML}{1D4C9A}
\definecolor{tintA}{HTML}{DBE5F6}
\definecolor{tintB}{HTML}{F4E4C6}
\definecolor{tintO}{HTML}{E6E6EA}
\definecolor{hair}{HTML}{A9ABB5}
\definecolor{old}{HTML}{B4B6C0}
\usepackage[colorlinks,allcolors=accent]{hyperref}

\newtheorem{theorem}{Theorem}
\newtheorem{lemma}{Lemma}
\newtheorem{corollary}[lemma]{Corollary}
\theoremstyle{definition}
\newtheorem{definition}{Definition}
\newtheorem*{example}{Example}
\theoremstyle{remark}
\newtheorem{remark}{Remark}

% A list whose items carry their own labels, such as (a) or (W1). The argument is the widest label.
\newenvironment{marked}[1]{%
  \begin{list}{}{\settowidth{\labelwidth}{\textup{#1}}%
    \setlength{\leftmargin}{\labelwidth}\addtolength{\leftmargin}{\labelsep}%
    \addtolength{\leftmargin}{1.3em}%
    \setlength{\itemsep}{1pt}\setlength{\parsep}{0pt}\setlength{\topsep}{3pt}%
    \setlength{\partopsep}{0pt}\renewcommand{\makelabel}[1]{\textup{##1}\hfil}}}%
  {\end{list}}

% Put at the start of a statement (theorem, lemma, corollary, definition, example): no page break
% inside it. \@endparpenalty is left alone, so that a page can still end right after the statement.
\makeatletter
\newcommand{\unbroken}{\interlinepenalty\@M \@beginparpenalty\@M \@itempenalty\@M
  \predisplaypenalty\@M \postdisplaypenalty\@M}
\makeatother

\title{The Falafel's walk on the square spiral never~ends}
\author{Martin Norbäck Olivers}
\date{Draft, 25 November 2026}

\begin{document}

\begin{abstract}
The Falafel is the chess piece that moves one or two squares diagonally. Put it on square 1 of the
square spiral and let it always move to the lowest-numbered square that it has not yet visited. The
numbers of the squares it visits form sequence A399160 of the OEIS. We prove that the Falafel is
never trapped and that it visits every square it can reach, so the sequence is a permutation of the
odd positive integers. After its first 80 squares the walk runs in laps around the origin, and we
give every lap in closed form. The proof is a case analysis in linear integer arithmetic. It has
been formalised in Lean~4 with Mathlib and is checked by the Lean kernel.
\end{abstract}

\maketitle

\section{Introduction}

Number the squares of an infinite chessboard along the square spiral: square 1 is \((0,0)\), square
2 is \((1,0)\), square 3 is \((1,1)\), and the numbering goes on counter-clockwise (Figure~1). Put a
chess piece on square 1 and let it move by one rule: it always moves to the lowest-numbered square
that it has not visited before. For the knight this walk ends. At step 2016 the knight stands on
square 2084, and every square one knight move away has been visited~\cite{ref2}.

\begin{figure}[htbp]
\centering
\begin{tikzpicture}[x=8.4mm,y=8.4mm]
\fill[tintA] (-3,3) rectangle +(1,1);
\fill[tintA] (-1,3) rectangle +(1,1);
\fill[tintA] (1,3) rectangle +(1,1);
\fill[tintA] (3,3) rectangle +(1,1);
\fill[tintA] (-2,2) rectangle +(1,1);
\fill[tintA] (0,2) rectangle +(1,1);
\fill[tintA] (2,2) rectangle +(1,1);
\fill[tintA] (-3,1) rectangle +(1,1);
\fill[tintA] (-1,1) rectangle +(1,1);
\fill[tintA] (1,1) rectangle +(1,1);
\fill[tintA] (3,1) rectangle +(1,1);
\fill[tintA] (-2,0) rectangle +(1,1);
\fill[tintA] (0,0) rectangle +(1,1);
\fill[tintA] (2,0) rectangle +(1,1);
\fill[tintA] (-3,-1) rectangle +(1,1);
\fill[tintA] (-1,-1) rectangle +(1,1);
\fill[tintA] (1,-1) rectangle +(1,1);
\fill[tintA] (3,-1) rectangle +(1,1);
\fill[tintA] (-2,-2) rectangle +(1,1);
\fill[tintA] (0,-2) rectangle +(1,1);
\fill[tintA] (2,-2) rectangle +(1,1);
\fill[tintA] (-3,-3) rectangle +(1,1);
\fill[tintA] (-1,-3) rectangle +(1,1);
\fill[tintA] (1,-3) rectangle +(1,1);
\fill[tintA] (3,-3) rectangle +(1,1);
\draw[hair,step=1] (-3,-3) grid (4,4);
\node at (-2.5,3.5) {37};
\node at (-1.5,3.5) {36};
\node at (-0.5,3.5) {35};
\node at (0.5,3.5) {34};
\node at (1.5,3.5) {33};
\node at (2.5,3.5) {32};
\node at (3.5,3.5) {31};
\node at (-2.5,2.5) {38};
\node at (-1.5,2.5) {17};
\node at (-0.5,2.5) {16};
\node at (0.5,2.5) {15};
\node at (1.5,2.5) {14};
\node at (2.5,2.5) {13};
\node at (3.5,2.5) {30};
\node at (-2.5,1.5) {39};
\node at (-1.5,1.5) {18};
\node at (-0.5,1.5) {5};
\node at (0.5,1.5) {4};
\node at (1.5,1.5) {3};
\node at (2.5,1.5) {12};
\node at (3.5,1.5) {29};
\node at (-2.5,0.5) {40};
\node at (-1.5,0.5) {19};
\node at (-0.5,0.5) {6};
\node at (0.5,0.5) {1};
\node at (1.5,0.5) {2};
\node at (2.5,0.5) {11};
\node at (3.5,0.5) {28};
\node at (-2.5,-0.5) {41};
\node at (-1.5,-0.5) {20};
\node at (-0.5,-0.5) {7};
\node at (0.5,-0.5) {8};
\node at (1.5,-0.5) {9};
\node at (2.5,-0.5) {10};
\node at (3.5,-0.5) {27};
\node at (-2.5,-1.5) {42};
\node at (-1.5,-1.5) {21};
\node at (-0.5,-1.5) {22};
\node at (0.5,-1.5) {23};
\node at (1.5,-1.5) {24};
\node at (2.5,-1.5) {25};
\node at (3.5,-1.5) {26};
\node at (-2.5,-2.5) {43};
\node at (-1.5,-2.5) {44};
\node at (-0.5,-2.5) {45};
\node at (0.5,-2.5) {46};
\node at (1.5,-2.5) {47};
\node at (2.5,-2.5) {48};
\node at (3.5,-2.5) {49};
\draw[line width=0.9pt] (-1.92,-1.92) rectangle +(0.84,0.84);
\draw[line width=0.9pt] (-1.92,2.08) rectangle +(0.84,0.84);
\draw[line width=0.9pt] (-0.92,-0.92) rectangle +(0.84,0.84);
\draw[line width=0.9pt] (-0.92,1.08) rectangle +(0.84,0.84);
\draw[line width=0.9pt] (1.08,-0.92) rectangle +(0.84,0.84);
\draw[line width=0.9pt] (1.08,1.08) rectangle +(0.84,0.84);
\draw[line width=0.9pt] (2.08,-1.92) rectangle +(0.84,0.84);
\draw[line width=0.9pt] (2.08,2.08) rectangle +(0.84,0.84);
\draw[accent,line width=1.3pt] (0.08,0.08) rectangle +(0.84,0.84);
\end{tikzpicture}
\caption{The square spiral. The Falafel stands only on the tinted squares, which are those with \(x + y\)
even. From square 1 it can move to the eight outlined squares 3, 5, 7, 9, 13, 17, 21 and 25.}
\end{figure}

Here the piece is the Falafel. It is the compound of the Fers, which moves one square diagonally,
and the Alfil, which leaps two squares diagonally. The numbers of the squares it visits are sequence
A399160 of the OEIS~\cite{ref1}, contributed by J.~Ahlström:
\[ 1,\ 3,\ 7,\ 19,\ 5,\ 9,\ 11,\ 15,\ 33,\ 13,\ 29,\ 53,\ \dots \]
On 23 November 2026 the entry did not say whether the sequence is finite. We prove that it is
infinite, and we determine it completely.

Write \(\sigma(x,y)\) for the number of the square \((x,y)\); the formula is (1) in Section~2. By
construction, every positive integer is the number of exactly one square. Two squares \((x,y)\) and
\((x',y')\) are \emph{one move apart} if \(|x'-x| = |y'-y|\) and this common value is 1 or 2.

\begin{definition}
\unbroken
A \emph{Falafel walk} is an infinite sequence \(a_0, a_1, a_2, \dots\) of squares such that

\begin{marked}{(W1)}
\item[(W1)] \(a_0 = (0,0)\);
\item[(W2)] \(a_m \ne a_n\) whenever \(m \ne n\);
\item[(W3)] \(a_{n+1}\) is one move from \(a_n\), for every \(n\);
\item[(W4)] for every \(n\), every square \(q\) that is one move from \(a_n\) and has \(\sigma(q) <
\sigma(a_{n+1})\) is one of \(a_0, \dots, a_n\).
\end{marked}
\end{definition}

\noindent
Conditions (W2) to (W4) say that \(a_{n+1}\) is the lowest-numbered unvisited square one move from
\(a_n\). A trapped piece produces only a finite sequence, so a Falafel walk exists exactly if the
Falafel is never trapped. The \(n\)-th term of A399160 is \(\sigma(a_{n-1})\).

\begin{remark}[only one walk]
Different squares have different numbers, so the rule allows at most one move: \(a_0, \dots, a_n\)
determine \(a_{n+1}\). So there is at most one Falafel walk.
\end{remark}

\begin{theorem}
\unbroken
\leavevmode
\begin{marked}{(a)}
\item[(a)] There is a Falafel walk.
\item[(b)] It visits every square \((x,y)\) with \(x+y\) even, and no other square.
\end{marked}
\end{theorem}

\noindent
The squares with \(x+y\) even are the squares with odd numbers. So A399160 is a permutation of the
odd positive integers (Corollary~8 in Section~6).

Theorem 2 in Section~3 describes the walk completely. Its first 80 squares are irregular. After them
the walk runs in laps \(R = 4, 5, 6, \dots\) around the origin, and one closed formula gives the
squares of every lap. Corollary~9 in Section~6 turns this into a formula for the \(n\)-th term of
the sequence.

The proof goes as follows. Every square gets a \emph{stamp}: the number of its lap, and its index
within the lap. Four lemmas about stamps carry the proof (Section~4). Each is a statement of linear
integer arithmetic and is proved by a case analysis, which we do not print. The whole proof has been
formalised in Lean~4~\cite{ref3} with the library Mathlib~\cite{ref4}, and the Lean kernel has
checked it. Sections 5 and 6 derive the theorems and Corollaries 8 and 9 from the lemmas by hand.
Section~7 says what was formalised, and which few lines a reader has to trust.

\section{The spiral in diagonal coordinates}

Let \(r = \max(|x|,|y|)\) be the \emph{ring} of the square \((x,y)\). The spiral number is
\begin{equation}
\sigma(x,y) \;=\; (2r-1)^2 + \begin{cases}
 y + r & \text{if } x = r \text{ and } y > -r, \\
 2r + (r - x) & \text{if } y = r, \\
 4r + (r - y) & \text{if } x = -r, \\
 6r + (x + r) & \text{otherwise,}
\end{cases}
\end{equation}
where the first case that applies is taken. Ring 0 is the square \((0,0)\) with number 1. Ring \(r
\ge 1\) holds the numbers from \((2r-1)^2 + 1\) to \((2r+1)^2\).

A move changes \(x\) by \(\pm 1\) or \(\pm 2\) and \(y\) by the same amount up to sign. So a move
does not change the parity of \(x+y\), and a Falafel walk, which starts at \((0,0)\), stays on the
squares with \(x+y\) even.

\begin{definition}[diagonal coordinates]
\unbroken
The \emph{diagonal coordinates} of a square \((x,y)\) with \(x+y\) even are the integers
\[ u = \frac{x+y}{2}, \qquad v = \frac{x-y}{2}, \qquad\text{so that}\qquad (x,y) = (u+v,\ u-v). \]
Every pair \((u,v)\) of integers belongs to exactly one such square. A \emph{step} is one of the
pairs \((\pm 1, 0)\), \((0, \pm 1)\), \((\pm 2, 0)\), \((0, \pm 2)\).
\end{definition}

\noindent
From now on, except in Section~6, a \emph{square} is given by its diagonal coordinates: it is a pair
\(s = (u,v)\) of integers, and \(\sigma(s)\) means \(\sigma(u+v,\,u-v)\). Two squares are one move
apart exactly if their difference is a step. So in diagonal coordinates the Falafel is a rook that
moves one or two squares.

\begin{definition}[ring and position]
\unbroken
For a square \((u,v)\) put \(\rho(u,v) = |u| + |v|\) and
\[ \pi(u,v) \;=\; \begin{cases} 2u & \text{if } u > 0,\ v \ge 0 \\ 2u - 4v & \text{if } u > 0,\ v < 0 \\ -6u - 4v & \text{if } u < 0,\ v \le 0 \\ 8v - 6u & \text{if } u < 0,\ v > 0 \\ -4v & \text{if } u = 0,\ v < 0 \\ 8v & \text{if } u = 0,\ v \ge 0 \end{cases} \]
\end{definition}

\begin{lemma}
\unbroken
Let \(s\) and \(s'\) be squares.

\begin{marked}{(iii)}
\item[(i)] \(\sigma(s) = (2\rho(s)-1)^2 + \pi(s)\).
\item[(ii)] \(0 \le \pi(s) \le 8\rho(s)\), and \(\pi(s) \ge 1\) if \(\rho(s) \ge 1\).
\item[(iii)] If \(\rho(s) < \rho(s')\), or \(\rho(s) = \rho(s')\) and \(\pi(s) \le \pi(s')\), then \(\sigma(s)
\le \sigma(s')\).
\end{marked}
\end{lemma}

\begin{proof}
Since \(\max(|u+v|,|u-v|) = |u|+|v|\), the number \(\rho(s)\) is the ring of \(s\). The cases of (1)
are the four sides of the ring. On each side, putting \(x = u+v\) and \(y = u-v\) into (1) gives the
stated value of \(\pi\). This is (i), and (ii) can be read off from the cases of \(\pi\). For (iii),
let \(\rho(s) < \rho(s')\). Then by (i) and (ii)
\[ \sigma(s) \;\le\; (2\rho(s)-1)^2 + 8\rho(s) \;=\; (2\rho(s)+1)^2 \;\le\; (2\rho(s')-1)^2 \;<\; \sigma(s'). \]
If the rings are equal, (iii) is (i).
\end{proof}

\section{The walk}

\begin{definition}[laps]
\unbroken
For an integer \(R \ge 4\) put
\[ \mathrm{lo}(R) = \begin{cases} -R - 1 & \text{if } R \text{ even} \\ -R - 3 & \text{if } R \text{ odd} \end{cases} \qquad \mathrm{hi}(R) = \begin{cases} (25R - 4)/2 & \text{if } R \text{ even} \\ (25R + 1)/2 & \text{if } R \text{ odd} \end{cases} \]
\emph{Lap} \(R\) is the sequence of squares \(P_R(j)\) for \(\mathrm{lo}(R) \le j <
\mathrm{hi}(R)\), where \(P_R(j)\) is given by Table~1.
\end{definition}

\noindent
A lap consists of nine pieces, A to I; Figure~2 shows lap 6. Piece A runs along the row \(v = R\)
(for odd \(R\) it starts one square below that row). Pieces B, D, F and H are staircases along the
four sides of a diamond, each one ring further out than the one before, and C, E and G are the
corners between them. Piece I comes back in along a slanted strip and ends one step from the start
of lap \(R+1\). The indices are counted from the start of piece B, so the indices of piece A are
negative.

\begin{figure}[htbp]
\centering
\begin{tikzpicture}[x=3.3mm,y=3.3mm]
\draw[hair] (-13,0) -- (15,0) node[below,text=black,font=\small] {$u$};
\draw[hair] (0,-13) -- (0,9) node[right,text=black,font=\small] {$v$};
\fill[old] (0,0) circle[radius=0.11];
\fill[old] (1,0) circle[radius=0.11];
\fill[old] (-1,0) circle[radius=0.11];
\fill[old] (-1,-1) circle[radius=0.11];
\fill[old] (0,-1) circle[radius=0.11];
\fill[old] (0,1) circle[radius=0.11];
\fill[old] (1,1) circle[radius=0.11];
\fill[old] (1,-1) circle[radius=0.11];
\fill[old] (2,-1) circle[radius=0.11];
\fill[old] (2,0) circle[radius=0.11];
\fill[old] (2,1) circle[radius=0.11];
\fill[old] (2,2) circle[radius=0.11];
\fill[old] (0,2) circle[radius=0.11];
\fill[old] (1,2) circle[radius=0.11];
\fill[old] (-1,2) circle[radius=0.11];
\fill[old] (-1,1) circle[radius=0.11];
\fill[old] (-2,1) circle[radius=0.11];
\fill[old] (-2,0) circle[radius=0.11];
\fill[old] (-2,-1) circle[radius=0.11];
\fill[old] (-2,-2) circle[radius=0.11];
\fill[old] (0,-2) circle[radius=0.11];
\fill[old] (1,-2) circle[radius=0.11];
\fill[old] (-1,-2) circle[radius=0.11];
\fill[old] (-1,-3) circle[radius=0.11];
\fill[old] (0,-3) circle[radius=0.11];
\fill[old] (1,-3) circle[radius=0.11];
\fill[old] (2,-3) circle[radius=0.11];
\fill[old] (2,-2) circle[radius=0.11];
\fill[old] (3,-2) circle[radius=0.11];
\fill[old] (3,0) circle[radius=0.11];
\fill[old] (3,1) circle[radius=0.11];
\fill[old] (3,-1) circle[radius=0.11];
\fill[old] (4,-1) circle[radius=0.11];
\fill[old] (4,0) circle[radius=0.11];
\fill[old] (4,1) circle[radius=0.11];
\fill[old] (4,2) circle[radius=0.11];
\fill[old] (3,2) circle[radius=0.11];
\fill[old] (3,3) circle[radius=0.11];
\fill[old] (1,3) circle[radius=0.11];
\fill[old] (0,3) circle[radius=0.11];
\fill[old] (-1,3) circle[radius=0.11];
\fill[old] (-2,3) circle[radius=0.11];
\fill[old] (-2,2) circle[radius=0.11];
\fill[old] (-3,2) circle[radius=0.11];
\fill[old] (-3,0) circle[radius=0.11];
\fill[old] (-3,-1) circle[radius=0.11];
\fill[old] (-3,1) circle[radius=0.11];
\fill[old] (-4,1) circle[radius=0.11];
\fill[old] (-4,0) circle[radius=0.11];
\fill[old] (-4,-1) circle[radius=0.11];
\fill[old] (-4,-2) circle[radius=0.11];
\fill[old] (-3,-2) circle[radius=0.11];
\fill[old] (-3,-3) circle[radius=0.11];
\fill[old] (-2,-3) circle[radius=0.11];
\fill[old] (-2,-4) circle[radius=0.11];
\fill[old] (0,-4) circle[radius=0.11];
\fill[old] (1,-4) circle[radius=0.11];
\fill[old] (-1,-4) circle[radius=0.11];
\fill[old] (-1,-5) circle[radius=0.11];
\fill[old] (0,-5) circle[radius=0.11];
\fill[old] (1,-5) circle[radius=0.11];
\fill[old] (2,-5) circle[radius=0.11];
\fill[old] (2,-4) circle[radius=0.11];
\fill[old] (3,-4) circle[radius=0.11];
\fill[old] (3,-3) circle[radius=0.11];
\fill[old] (4,-3) circle[radius=0.11];
\fill[old] (4,-2) circle[radius=0.11];
\fill[old] (5,-2) circle[radius=0.11];
\fill[old] (5,0) circle[radius=0.11];
\fill[old] (5,1) circle[radius=0.11];
\fill[old] (5,-1) circle[radius=0.11];
\fill[old] (6,-1) circle[radius=0.11];
\fill[old] (6,0) circle[radius=0.11];
\fill[old] (6,1) circle[radius=0.11];
\fill[old] (6,2) circle[radius=0.11];
\fill[old] (5,2) circle[radius=0.11];
\fill[old] (5,3) circle[radius=0.11];
\fill[old] (4,3) circle[radius=0.11];
\fill[old] (2,3) circle[radius=0.11];
\fill[old] (2,4) circle[radius=0.11];
\fill[old] (0,4) circle[radius=0.11];
\fill[old] (1,4) circle[radius=0.11];
\fill[old] (-1,4) circle[radius=0.11];
\fill[old] (-2,4) circle[radius=0.11];
\fill[old] (-3,4) circle[radius=0.11];
\fill[old] (-3,3) circle[radius=0.11];
\fill[old] (-4,3) circle[radius=0.11];
\fill[old] (-4,2) circle[radius=0.11];
\fill[old] (-5,2) circle[radius=0.11];
\fill[old] (-5,0) circle[radius=0.11];
\fill[old] (-5,-1) circle[radius=0.11];
\fill[old] (-5,1) circle[radius=0.11];
\fill[old] (-6,1) circle[radius=0.11];
\fill[old] (-6,0) circle[radius=0.11];
\fill[old] (-6,-1) circle[radius=0.11];
\fill[old] (-6,-2) circle[radius=0.11];
\fill[old] (-5,-2) circle[radius=0.11];
\fill[old] (-5,-3) circle[radius=0.11];
\fill[old] (-4,-3) circle[radius=0.11];
\fill[old] (-4,-4) circle[radius=0.11];
\fill[old] (-3,-4) circle[radius=0.11];
\fill[old] (-3,-5) circle[radius=0.11];
\fill[old] (-2,-5) circle[radius=0.11];
\fill[old] (-2,-6) circle[radius=0.11];
\fill[old] (0,-6) circle[radius=0.11];
\fill[old] (1,-6) circle[radius=0.11];
\fill[old] (-1,-6) circle[radius=0.11];
\fill[old] (-1,-7) circle[radius=0.11];
\fill[old] (0,-7) circle[radius=0.11];
\fill[old] (1,-7) circle[radius=0.11];
\fill[old] (2,-7) circle[radius=0.11];
\fill[old] (2,-6) circle[radius=0.11];
\fill[old] (3,-6) circle[radius=0.11];
\fill[old] (3,-5) circle[radius=0.11];
\fill[old] (4,-5) circle[radius=0.11];
\fill[old] (4,-4) circle[radius=0.11];
\fill[old] (5,-4) circle[radius=0.11];
\fill[old] (5,-3) circle[radius=0.11];
\fill[old] (6,-3) circle[radius=0.11];
\fill[old] (6,-2) circle[radius=0.11];
\fill[old] (7,-2) circle[radius=0.11];
\fill[old] (7,0) circle[radius=0.11];
\fill[old] (7,1) circle[radius=0.11];
\fill[old] (7,-1) circle[radius=0.11];
\fill[old] (8,-1) circle[radius=0.11];
\fill[old] (8,0) circle[radius=0.11];
\fill[old] (8,1) circle[radius=0.11];
\fill[old] (8,2) circle[radius=0.11];
\fill[old] (7,2) circle[radius=0.11];
\fill[old] (7,3) circle[radius=0.11];
\fill[old] (6,3) circle[radius=0.11];
\fill[old] (6,4) circle[radius=0.11];
\fill[old] (4,4) circle[radius=0.11];
\fill[old] (3,4) circle[radius=0.11];
\fill[old] (3,5) circle[radius=0.11];
\fill[old] (1,5) circle[radius=0.11];
\fill[old] (0,5) circle[radius=0.11];
\fill[old] (-1,5) circle[radius=0.11];
\fill[old] (-2,5) circle[radius=0.11];
\fill[old] (-3,5) circle[radius=0.11];
\fill[old] (-4,5) circle[radius=0.11];
\fill[old] (-4,4) circle[radius=0.11];
\fill[old] (-5,4) circle[radius=0.11];
\fill[old] (-5,3) circle[radius=0.11];
\fill[old] (-6,3) circle[radius=0.11];
\fill[old] (-6,2) circle[radius=0.11];
\fill[old] (-7,2) circle[radius=0.11];
\fill[old] (-7,0) circle[radius=0.11];
\fill[old] (-7,-1) circle[radius=0.11];
\fill[old] (-7,1) circle[radius=0.11];
\fill[old] (-8,1) circle[radius=0.11];
\fill[old] (-8,0) circle[radius=0.11];
\fill[old] (-8,-1) circle[radius=0.11];
\fill[old] (-8,-2) circle[radius=0.11];
\fill[old] (-7,-2) circle[radius=0.11];
\fill[old] (-7,-3) circle[radius=0.11];
\fill[old] (-6,-3) circle[radius=0.11];
\fill[old] (-6,-4) circle[radius=0.11];
\fill[old] (-5,-4) circle[radius=0.11];
\fill[old] (-5,-5) circle[radius=0.11];
\fill[old] (-4,-5) circle[radius=0.11];
\fill[old] (-4,-6) circle[radius=0.11];
\fill[old] (-3,-6) circle[radius=0.11];
\fill[old] (-3,-7) circle[radius=0.11];
\fill[old] (-2,-7) circle[radius=0.11];
\fill[old] (-2,-8) circle[radius=0.11];
\fill[old] (0,-8) circle[radius=0.11];
\fill[old] (1,-8) circle[radius=0.11];
\fill[old] (-1,-8) circle[radius=0.11];
\fill[old] (-1,-9) circle[radius=0.11];
\fill[old] (0,-9) circle[radius=0.11];
\fill[old] (1,-9) circle[radius=0.11];
\fill[old] (2,-9) circle[radius=0.11];
\fill[old] (2,-8) circle[radius=0.11];
\fill[old] (3,-8) circle[radius=0.11];
\fill[old] (3,-7) circle[radius=0.11];
\fill[old] (4,-7) circle[radius=0.11];
\fill[old] (4,-6) circle[radius=0.11];
\fill[old] (5,-6) circle[radius=0.11];
\fill[old] (5,-5) circle[radius=0.11];
\fill[old] (6,-5) circle[radius=0.11];
\fill[old] (6,-4) circle[radius=0.11];
\fill[old] (7,-4) circle[radius=0.11];
\fill[old] (7,-3) circle[radius=0.11];
\fill[old] (8,-3) circle[radius=0.11];
\fill[old] (8,-2) circle[radius=0.11];
\fill[old] (9,-2) circle[radius=0.11];
\fill[old] (9,0) circle[radius=0.11];
\fill[old] (9,1) circle[radius=0.11];
\fill[old] (9,-1) circle[radius=0.11];
\fill[old] (10,-1) circle[radius=0.11];
\fill[old] (10,0) circle[radius=0.11];
\fill[old] (10,1) circle[radius=0.11];
\fill[old] (10,2) circle[radius=0.11];
\fill[old] (9,2) circle[radius=0.11];
\fill[old] (9,3) circle[radius=0.11];
\fill[old] (8,3) circle[radius=0.11];
\fill[old] (8,4) circle[radius=0.11];
\fill[old] (7,4) circle[radius=0.11];
\fill[old] (5,4) circle[radius=0.11];
\fill[old] (5,5) circle[radius=0.11];
\fill[old] (4,5) circle[radius=0.11];
\fill[old] (2,5) circle[radius=0.11];
\fill[old] (2,6) circle[radius=0.11];
\draw[accent,line width=0.7pt,line join=round]
  (0,6) -- (1,6) -- (-1,6) -- (-2,6) -- (-3,6) -- (-4,6) -- (-5,6) -- (-5,5) -- (-6,5) -- (-6,4) --
  (-7,4) -- (-7,3) -- (-8,3) -- (-8,2) -- (-9,2) -- (-9,0) -- (-9,-1) -- (-9,1) -- (-10,1) -- (-10,0) --
  (-10,-1) -- (-10,-2) -- (-9,-2) -- (-9,-3) -- (-8,-3) -- (-8,-4) -- (-7,-4) -- (-7,-5) -- (-6,-5) -- (-6,-6) --
  (-5,-6) -- (-5,-7) -- (-4,-7) -- (-4,-8) -- (-3,-8) -- (-3,-9) -- (-2,-9) -- (-2,-10) -- (0,-10) -- (1,-10) --
  (-1,-10) -- (-1,-11) -- (0,-11) -- (1,-11) -- (2,-11) -- (2,-10) -- (3,-10) -- (3,-9) -- (4,-9) -- (4,-8) --
  (5,-8) -- (5,-7) -- (6,-7) -- (6,-6) -- (7,-6) -- (7,-5) -- (8,-5) -- (8,-4) -- (9,-4) -- (9,-3) --
  (10,-3) -- (10,-2) -- (11,-2) -- (11,0) -- (11,1) -- (11,-1) -- (12,-1) -- (12,0) -- (12,1) -- (12,2) --
  (11,2) -- (11,3) -- (10,3) -- (10,4) -- (9,4) -- (9,5) -- (7,5) -- (6,5) -- (6,6) -- (4,6);
\fill[accent] (0,6) circle[radius=0.14];
\fill[accent] (1,6) circle[radius=0.14];
\fill[accent] (-1,6) circle[radius=0.14];
\fill[accent] (-2,6) circle[radius=0.14];
\fill[accent] (-3,6) circle[radius=0.14];
\fill[accent] (-4,6) circle[radius=0.14];
\fill[accent] (-5,6) circle[radius=0.14];
\fill[accent] (-5,5) circle[radius=0.14];
\fill[accent] (-6,5) circle[radius=0.14];
\fill[accent] (-6,4) circle[radius=0.14];
\fill[accent] (-7,4) circle[radius=0.14];
\fill[accent] (-7,3) circle[radius=0.14];
\fill[accent] (-8,3) circle[radius=0.14];
\fill[accent] (-8,2) circle[radius=0.14];
\fill[accent] (-9,2) circle[radius=0.14];
\fill[accent] (-9,0) circle[radius=0.14];
\fill[accent] (-9,-1) circle[radius=0.14];
\fill[accent] (-9,1) circle[radius=0.14];
\fill[accent] (-10,1) circle[radius=0.14];
\fill[accent] (-10,0) circle[radius=0.14];
\fill[accent] (-10,-1) circle[radius=0.14];
\fill[accent] (-10,-2) circle[radius=0.14];
\fill[accent] (-9,-2) circle[radius=0.14];
\fill[accent] (-9,-3) circle[radius=0.14];
\fill[accent] (-8,-3) circle[radius=0.14];
\fill[accent] (-8,-4) circle[radius=0.14];
\fill[accent] (-7,-4) circle[radius=0.14];
\fill[accent] (-7,-5) circle[radius=0.14];
\fill[accent] (-6,-5) circle[radius=0.14];
\fill[accent] (-6,-6) circle[radius=0.14];
\fill[accent] (-5,-6) circle[radius=0.14];
\fill[accent] (-5,-7) circle[radius=0.14];
\fill[accent] (-4,-7) circle[radius=0.14];
\fill[accent] (-4,-8) circle[radius=0.14];
\fill[accent] (-3,-8) circle[radius=0.14];
\fill[accent] (-3,-9) circle[radius=0.14];
\fill[accent] (-2,-9) circle[radius=0.14];
\fill[accent] (-2,-10) circle[radius=0.14];
\fill[accent] (0,-10) circle[radius=0.14];
\fill[accent] (1,-10) circle[radius=0.14];
\fill[accent] (-1,-10) circle[radius=0.14];
\fill[accent] (-1,-11) circle[radius=0.14];
\fill[accent] (0,-11) circle[radius=0.14];
\fill[accent] (1,-11) circle[radius=0.14];
\fill[accent] (2,-11) circle[radius=0.14];
\fill[accent] (2,-10) circle[radius=0.14];
\fill[accent] (3,-10) circle[radius=0.14];
\fill[accent] (3,-9) circle[radius=0.14];
\fill[accent] (4,-9) circle[radius=0.14];
\fill[accent] (4,-8) circle[radius=0.14];
\fill[accent] (5,-8) circle[radius=0.14];
\fill[accent] (5,-7) circle[radius=0.14];
\fill[accent] (6,-7) circle[radius=0.14];
\fill[accent] (6,-6) circle[radius=0.14];
\fill[accent] (7,-6) circle[radius=0.14];
\fill[accent] (7,-5) circle[radius=0.14];
\fill[accent] (8,-5) circle[radius=0.14];
\fill[accent] (8,-4) circle[radius=0.14];
\fill[accent] (9,-4) circle[radius=0.14];
\fill[accent] (9,-3) circle[radius=0.14];
\fill[accent] (10,-3) circle[radius=0.14];
\fill[accent] (10,-2) circle[radius=0.14];
\fill[accent] (11,-2) circle[radius=0.14];
\fill[accent] (11,0) circle[radius=0.14];
\fill[accent] (11,1) circle[radius=0.14];
\fill[accent] (11,-1) circle[radius=0.14];
\fill[accent] (12,-1) circle[radius=0.14];
\fill[accent] (12,0) circle[radius=0.14];
\fill[accent] (12,1) circle[radius=0.14];
\fill[accent] (12,2) circle[radius=0.14];
\fill[accent] (11,2) circle[radius=0.14];
\fill[accent] (11,3) circle[radius=0.14];
\fill[accent] (10,3) circle[radius=0.14];
\fill[accent] (10,4) circle[radius=0.14];
\fill[accent] (9,4) circle[radius=0.14];
\fill[accent] (9,5) circle[radius=0.14];
\fill[accent] (7,5) circle[radius=0.14];
\fill[accent] (6,5) circle[radius=0.14];
\fill[accent] (6,6) circle[radius=0.14];
\fill[accent] (4,6) circle[radius=0.14];
\fill[black] (0,6) circle[radius=0.27];
\draw[black,line width=0.9pt,fill=white] (4,6) circle[radius=0.25];
\node[font=\bfseries\itshape] at (-3.00,7.10) {A};
\node[font=\bfseries\itshape] at (-8.34,4.17) {B};
\node[font=\bfseries\itshape] at (-11.00,0.00) {C};
\node[font=\bfseries\itshape] at (-6.83,-6.83) {D};
\node[font=\bfseries\itshape] at (0.00,-12.00) {E};
\node[font=\bfseries\itshape] at (7.32,-7.32) {F};
\node[font=\bfseries\itshape] at (13.00,0.00) {G};
\node[font=\bfseries\itshape] at (11.71,3.74) {H};
\node[font=\bfseries\itshape] at (6.95,6.90) {I};
\end{tikzpicture}
\caption{Lap 6 in the coordinates \((u,v)\). It runs from \((0,6)\), the filled circle, to \((4,6)\), the
open circle. The letters are the pieces of Table~1. The grey dots are the squares visited before lap
6.}
\end{figure}

\begin{table}[htbp]
\caption{The square \(P_R(j)\) of lap \(R \ge 4\) at index \(j\). In piece I, \(h = \lfloor R/2 \rfloor\) and
\(k = \lfloor t/3 \rfloor\).}
\centering\footnotesize
\begin{tabular}{@{}ll>{\raggedright\arraybackslash}p{0.42\textwidth}@{}}
\toprule
\textit{Piece} & \textit{Index \(j\)} & \textit{Square \(P_R(j)\)} \\
\midrule
A. North row & even \(R\), \(j = -R - 1\) & \((0,\ R)\) \\
 & even \(R\), \(j = -R\) & \((1,\ R)\) \\
 & odd \(R\), \(j = -R - 3\) & \((3,\ R - 1)\) \\
 & odd \(R\), \(j = -R - 2\) & \((3,\ R)\) \\
 & all other \(j\) with \(\mathrm{lo}(R) \le j < 0\) & \((-R - j,\ R)\) \\
\addlinespace
B. North-west staircase & \(j = t\), \(0 \le t < 2R - 4\) & \((1 - R - \lfloor (t + 1)/2 \rfloor,\ R - 1 - \lfloor t/2 \rfloor)\) \\
\addlinespace
C. West corner & \(j = 2R - 4 + t\), \(t = 0, \dots, 5\) & \((3 - 2R,\,0),\) \((3 - 2R,\,-1),\) \((3 - 2R,\,1),\) \((2 - 2R,\,1),\) \((2 - 2R,\,0),\) \((2 - 2R,\,-1)\) \\
\addlinespace
D. South-west staircase & \(j = 2R + 2 + t\), \(0 \le t < 4R - 7\) & \((2 - 2R + \lfloor (t + 1)/2 \rfloor,\ -2 - \lfloor t/2 \rfloor)\) \\
\addlinespace
E. South corner & \(j = 6R - 5 + t\), \(t = 0, \dots, 5\) & \((0,\,2 - 2R),\) \((1,\,2 - 2R),\) \((-1,\,2 - 2R),\) \((-1,\,1 - 2R),\) \((0,\,1 - 2R),\) \((1,\,1 - 2R)\) \\
\addlinespace
F. South-east staircase & \(j = 6R + 1 + t\), \(0 \le t < 4R - 5\) & \((2 + \lfloor t/2 \rfloor,\ 1 - 2R + \lfloor (t + 1)/2 \rfloor)\) \\
\addlinespace
G. East corner & \(j = 10R - 4 + t\), \(t = 0, \dots, 5\) & \((2R - 1,\,0),\) \((2R - 1,\,1),\) \((2R - 1,\,-1),\) \((2R,\,-1),\) \((2R,\,0),\) \((2R,\,1)\) \\
\addlinespace
H. North-east staircase & \(j = 10R + 2 + t\), \(0 \le t < R + 1\) & \((2R - \lfloor (t + 1)/2 \rfloor,\ 2 + \lfloor t/2 \rfloor)\) \\
\addlinespace
I. Return & \(j = 11R + 3 + t < \mathrm{hi}(R)\), \(t \ge 0\) &  \\
 & even \(R\), \(t = 3k\) & \((3h - 3k - 2,\ h + k + 2)\) \\
 & even \(R\), \(t = 3k + 1\) & \((3h - 3k - 3,\ h + k + 2)\) \\
 & even \(R\), \(t = 3k + 2\) & \((3h - 3k - 3,\ h + k + 3)\) \\
 & odd \(R\), \(t = 3k\) & \((3h - 3k - 1,\ h + k + 2)\) \\
 & odd \(R\), \(t = 3k + 1\) & \((3h - 3k - 1,\ h + k + 3)\) \\
 & odd \(R\), \(t = 3k + 2\) & \((3h - 3k - 2,\ h + k + 3)\) \\
\bottomrule
\end{tabular}
\end{table}

Let \(T_0, T_1, \dots, T_{79}\) be the 80 squares of Table~4 in Appendix~A.

\begin{theorem}
\unbroken
Let \(W_0, W_1, W_2, \dots\) be the sequence that consists of \(T_0, \dots, T_{79}\), then lap 4,
then lap 5, and so on.

\begin{marked}{(a)}
\item[(a)] Every square \((u,v)\) of \(\mathbb{Z}^2\) occurs exactly once in this sequence.
\item[(b)] Write \(W_n = (u_n, v_n)\). Then \(a_n = (u_n + v_n,\ u_n - v_n)\) is the Falafel walk.
\end{marked}
\end{theorem}

\noindent
The number of squares of lap \(R\) is
\[ \mathrm{hi}(R) - \mathrm{lo}(R) = \begin{cases} (27R - 2)/2 & \text{if } R \text{ even} \\ (27R + 7)/2 & \text{if } R \text{ odd} \end{cases} \]
So lap 4 starts with \(W_{80}\) and lap 5 with \(W_{133}\). Corollary~9 gives the start of every
lap.

A program checked the following, and the description of the pieces above, for \(4 \le R \le 300\);
it is not part of the formal proof. Lap \(R\) lies in the rings \(R \le \rho \le 2R + 2\) and has
squares in both of these rings. So when lap \(R\) is complete, every ring up to \(R\) is fully
visited, ring \(R + 1\) is not, and no square beyond ring \(2R + 2\) has been visited. The outermost
visited ring is about twice the first incomplete one.

\section{Stamps}

Table~1 gives the square at a given index of a lap. The next definition goes the other way: it gives
the lap and the index of a given square (Lemma~3).

\begin{definition}[lap number and index]
\unbroken
The \emph{lap number} \(\lambda(u,v)\) of a square is given by Table~2. For a square with
\(\lambda(u,v) \ge 4\), its \emph{index} \(\kappa(u,v)\) is given by Table~5 in Appendix~A.
\end{definition}

\noindent
The ten regions of Table~2 and the centre are disjoint and cover the plane. Regions 1 to 8 hold the
pieces A to H, and region 10 holds piece I; region 9 holds the first two squares of an odd lap. A
value below 4 is not the number of a lap: by Lemma~2 it marks one of the first 80 squares. Figure~3
shows the lap numbers near the origin.

\begin{table}[htbp]
\caption{The lap number \(\lambda(u,v)\). In region 10, \(q = \lfloor (u + 3v)/6 \rfloor\) and \(m = (u + 3v)
\bmod 6\).}
\centering\small
\begin{tabular}{@{}llll@{}}
\toprule
\textit{Region} & \textit{Squares \((u, v)\)} & \textit{Case} & \(\lambda(u, v)\) \\
\midrule
1. North row & \(v \ge 2,\ u \le 1,\ u + v \ge 1\) &  & \(v\) \\
2. North-west staircase & \(v \ge 2,\ u + v \le 0\) &  & \(\lfloor (v - u)/2 \rfloor + 1\) \\
3. West corner & \(|v| \le 1,\ u \le -2\) &  & \(\lfloor (3 - u)/2 \rfloor\) \\
4. South-west staircase & \(v \le -2,\ u \le -2\) &  & \(\lfloor (1 - u - v)/2 \rfloor\) \\
5. South corner & \(v \le -2,\ |u| \le 1\) &  & \(\lfloor (2 - v)/2 \rfloor\) \\
6. South-east staircase & \(v \le -2,\ u \ge 2\) &  & \(\lfloor (u - v)/2 \rfloor\) \\
7. East corner & \(|v| \le 1,\ u \ge 2\) &  & \(\lfloor (u + 1)/2 \rfloor\) \\
8. North-east staircase & \(v \ge 2,\ u \ge 2,\ u - 3v \ge -6\) &  & \(\lfloor (u + v - 1)/2 \rfloor\) \\
9. Column 3 & \(u = 3,\ v \ge 4\) &  & \(2\lfloor v/2 \rfloor + 1\) \\
10. Return & \(v \ge 2,\ u \ge 2,\ u \ne 3,\ u - 3v \le -7\) & \(m = 0\) & \(2q - 2\) \\
 &  & \(1 \le m \le 2\) & \(2q - 1\) \\
 &  & \(3 \le m \le 4\) & \(2q\) \\
 &  & \(m = 5\) & \(2q + 1\) \\
Centre & \(|u| \le 1,\ |v| \le 1\) &  & \(0\) \\
\bottomrule
\end{tabular}
\end{table}

\begin{figure}[htbp]
\centering
\begin{tikzpicture}[x=6.4mm,y=6.4mm,every node/.style={font=\footnotesize,inner sep=0pt}]
\fill[tintA] (-9,9) rectangle +(1,1);
\fill[tintB] (-8,9) rectangle +(1,1);
\fill[tintB] (-7,9) rectangle +(1,1);
\fill[tintB] (-6,9) rectangle +(1,1);
\fill[tintB] (-5,9) rectangle +(1,1);
\fill[tintB] (-4,9) rectangle +(1,1);
\fill[tintB] (-3,9) rectangle +(1,1);
\fill[tintB] (-2,9) rectangle +(1,1);
\fill[tintB] (-1,9) rectangle +(1,1);
\fill[tintB] (0,9) rectangle +(1,1);
\fill[tintB] (1,9) rectangle +(1,1);
\fill[tintB] (2,9) rectangle +(1,1);
\fill[tintB] (3,9) rectangle +(1,1);
\fill[tintB] (4,9) rectangle +(1,1);
\fill[tintB] (5,9) rectangle +(1,1);
\fill[tintA] (6,9) rectangle +(1,1);
\fill[tintA] (7,9) rectangle +(1,1);
\fill[tintB] (8,9) rectangle +(1,1);
\fill[tintA] (9,9) rectangle +(1,1);
\fill[tintB] (-9,8) rectangle +(1,1);
\fill[tintB] (-8,8) rectangle +(1,1);
\fill[tintA] (-7,8) rectangle +(1,1);
\fill[tintA] (-6,8) rectangle +(1,1);
\fill[tintA] (-5,8) rectangle +(1,1);
\fill[tintA] (-4,8) rectangle +(1,1);
\fill[tintA] (-3,8) rectangle +(1,1);
\fill[tintA] (-2,8) rectangle +(1,1);
\fill[tintA] (-1,8) rectangle +(1,1);
\fill[tintA] (0,8) rectangle +(1,1);
\fill[tintA] (1,8) rectangle +(1,1);
\fill[tintB] (2,8) rectangle +(1,1);
\fill[tintB] (3,8) rectangle +(1,1);
\fill[tintA] (4,8) rectangle +(1,1);
\fill[tintB] (5,8) rectangle +(1,1);
\fill[tintA] (6,8) rectangle +(1,1);
\fill[tintB] (7,8) rectangle +(1,1);
\fill[tintB] (8,8) rectangle +(1,1);
\fill[tintA] (9,8) rectangle +(1,1);
\fill[tintB] (-9,7) rectangle +(1,1);
\fill[tintA] (-8,7) rectangle +(1,1);
\fill[tintA] (-7,7) rectangle +(1,1);
\fill[tintB] (-6,7) rectangle +(1,1);
\fill[tintB] (-5,7) rectangle +(1,1);
\fill[tintB] (-4,7) rectangle +(1,1);
\fill[tintB] (-3,7) rectangle +(1,1);
\fill[tintB] (-2,7) rectangle +(1,1);
\fill[tintB] (-1,7) rectangle +(1,1);
\fill[tintB] (0,7) rectangle +(1,1);
\fill[tintB] (1,7) rectangle +(1,1);
\fill[tintB] (2,7) rectangle +(1,1);
\fill[tintB] (3,7) rectangle +(1,1);
\fill[tintB] (4,7) rectangle +(1,1);
\fill[tintB] (5,7) rectangle +(1,1);
\fill[tintA] (6,7) rectangle +(1,1);
\fill[tintA] (7,7) rectangle +(1,1);
\fill[tintB] (8,7) rectangle +(1,1);
\fill[tintA] (9,7) rectangle +(1,1);
\fill[tintA] (-9,6) rectangle +(1,1);
\fill[tintA] (-8,6) rectangle +(1,1);
\fill[tintB] (-7,6) rectangle +(1,1);
\fill[tintB] (-6,6) rectangle +(1,1);
\fill[tintA] (-5,6) rectangle +(1,1);
\fill[tintA] (-4,6) rectangle +(1,1);
\fill[tintA] (-3,6) rectangle +(1,1);
\fill[tintA] (-2,6) rectangle +(1,1);
\fill[tintA] (-1,6) rectangle +(1,1);
\fill[tintA] (0,6) rectangle +(1,1);
\fill[tintA] (1,6) rectangle +(1,1);
\fill[tintB] (2,6) rectangle +(1,1);
\fill[tintB] (3,6) rectangle +(1,1);
\fill[tintA] (4,6) rectangle +(1,1);
\fill[tintB] (5,6) rectangle +(1,1);
\fill[tintA] (6,6) rectangle +(1,1);
\fill[tintB] (7,6) rectangle +(1,1);
\fill[tintB] (8,6) rectangle +(1,1);
\fill[tintA] (9,6) rectangle +(1,1);
\fill[tintA] (-9,5) rectangle +(1,1);
\fill[tintB] (-8,5) rectangle +(1,1);
\fill[tintB] (-7,5) rectangle +(1,1);
\fill[tintA] (-6,5) rectangle +(1,1);
\fill[tintA] (-5,5) rectangle +(1,1);
\fill[tintB] (-4,5) rectangle +(1,1);
\fill[tintB] (-3,5) rectangle +(1,1);
\fill[tintB] (-2,5) rectangle +(1,1);
\fill[tintB] (-1,5) rectangle +(1,1);
\fill[tintB] (0,5) rectangle +(1,1);
\fill[tintB] (1,5) rectangle +(1,1);
\fill[tintB] (2,5) rectangle +(1,1);
\fill[tintB] (3,5) rectangle +(1,1);
\fill[tintB] (4,5) rectangle +(1,1);
\fill[tintB] (5,5) rectangle +(1,1);
\fill[tintA] (6,5) rectangle +(1,1);
\fill[tintA] (7,5) rectangle +(1,1);
\fill[tintB] (8,5) rectangle +(1,1);
\fill[tintA] (9,5) rectangle +(1,1);
\fill[tintB] (-9,4) rectangle +(1,1);
\fill[tintB] (-8,4) rectangle +(1,1);
\fill[tintA] (-7,4) rectangle +(1,1);
\fill[tintA] (-6,4) rectangle +(1,1);
\fill[tintB] (-5,4) rectangle +(1,1);
\fill[tintB] (-4,4) rectangle +(1,1);
\fill[tintA] (-3,4) rectangle +(1,1);
\fill[tintA] (-2,4) rectangle +(1,1);
\fill[tintA] (-1,4) rectangle +(1,1);
\fill[tintA] (0,4) rectangle +(1,1);
\fill[tintA] (1,4) rectangle +(1,1);
\fill[tintO] (2,4) rectangle +(1,1);
\fill[tintB] (3,4) rectangle +(1,1);
\fill[tintA] (4,4) rectangle +(1,1);
\fill[tintB] (5,4) rectangle +(1,1);
\fill[tintA] (6,4) rectangle +(1,1);
\fill[tintB] (7,4) rectangle +(1,1);
\fill[tintB] (8,4) rectangle +(1,1);
\fill[tintA] (9,4) rectangle +(1,1);
\fill[tintB] (-9,3) rectangle +(1,1);
\fill[tintA] (-8,3) rectangle +(1,1);
\fill[tintA] (-7,3) rectangle +(1,1);
\fill[tintB] (-6,3) rectangle +(1,1);
\fill[tintB] (-5,3) rectangle +(1,1);
\fill[tintA] (-4,3) rectangle +(1,1);
\fill[tintA] (-3,3) rectangle +(1,1);
\fill[tintO] (-2,3) rectangle +(1,1);
\fill[tintO] (-1,3) rectangle +(1,1);
\fill[tintO] (0,3) rectangle +(1,1);
\fill[tintO] (1,3) rectangle +(1,1);
\fill[tintO] (2,3) rectangle +(1,1);
\fill[tintO] (3,3) rectangle +(1,1);
\fill[tintO] (4,3) rectangle +(1,1);
\fill[tintO] (5,3) rectangle +(1,1);
\fill[tintA] (6,3) rectangle +(1,1);
\fill[tintA] (7,3) rectangle +(1,1);
\fill[tintB] (8,3) rectangle +(1,1);
\fill[tintB] (9,3) rectangle +(1,1);
\fill[tintA] (-9,2) rectangle +(1,1);
\fill[tintA] (-8,2) rectangle +(1,1);
\fill[tintB] (-7,2) rectangle +(1,1);
\fill[tintB] (-6,2) rectangle +(1,1);
\fill[tintA] (-5,2) rectangle +(1,1);
\fill[tintA] (-4,2) rectangle +(1,1);
\fill[tintO] (-3,2) rectangle +(1,1);
\fill[tintO] (-2,2) rectangle +(1,1);
\fill[tintO] (-1,2) rectangle +(1,1);
\fill[tintO] (0,2) rectangle +(1,1);
\fill[tintO] (1,2) rectangle +(1,1);
\fill[tintO] (2,2) rectangle +(1,1);
\fill[tintO] (3,2) rectangle +(1,1);
\fill[tintO] (4,2) rectangle +(1,1);
\fill[tintO] (5,2) rectangle +(1,1);
\fill[tintO] (6,2) rectangle +(1,1);
\fill[tintA] (7,2) rectangle +(1,1);
\fill[tintA] (8,2) rectangle +(1,1);
\fill[tintB] (9,2) rectangle +(1,1);
\fill[tintA] (-9,1) rectangle +(1,1);
\fill[tintB] (-8,1) rectangle +(1,1);
\fill[tintB] (-7,1) rectangle +(1,1);
\fill[tintA] (-6,1) rectangle +(1,1);
\fill[tintA] (-5,1) rectangle +(1,1);
\fill[tintO] (-4,1) rectangle +(1,1);
\fill[tintO] (-3,1) rectangle +(1,1);
\fill[tintO] (-2,1) rectangle +(1,1);
\fill[tintO] (-1,1) rectangle +(1,1);
\fill[tintO] (0,1) rectangle +(1,1);
\fill[tintO] (1,1) rectangle +(1,1);
\fill[tintO] (2,1) rectangle +(1,1);
\fill[tintO] (3,1) rectangle +(1,1);
\fill[tintO] (4,1) rectangle +(1,1);
\fill[tintO] (5,1) rectangle +(1,1);
\fill[tintO] (6,1) rectangle +(1,1);
\fill[tintA] (7,1) rectangle +(1,1);
\fill[tintA] (8,1) rectangle +(1,1);
\fill[tintB] (9,1) rectangle +(1,1);
\fill[tintA] (-9,0) rectangle +(1,1);
\fill[tintB] (-8,0) rectangle +(1,1);
\fill[tintB] (-7,0) rectangle +(1,1);
\fill[tintA] (-6,0) rectangle +(1,1);
\fill[tintA] (-5,0) rectangle +(1,1);
\fill[tintO] (-4,0) rectangle +(1,1);
\fill[tintO] (-3,0) rectangle +(1,1);
\fill[tintO] (-2,0) rectangle +(1,1);
\fill[tintO] (-1,0) rectangle +(1,1);
\fill[tintO] (0,0) rectangle +(1,1);
\fill[tintO] (1,0) rectangle +(1,1);
\fill[tintO] (2,0) rectangle +(1,1);
\fill[tintO] (3,0) rectangle +(1,1);
\fill[tintO] (4,0) rectangle +(1,1);
\fill[tintO] (5,0) rectangle +(1,1);
\fill[tintO] (6,0) rectangle +(1,1);
\fill[tintA] (7,0) rectangle +(1,1);
\fill[tintA] (8,0) rectangle +(1,1);
\fill[tintB] (9,0) rectangle +(1,1);
\fill[tintA] (-9,-1) rectangle +(1,1);
\fill[tintB] (-8,-1) rectangle +(1,1);
\fill[tintB] (-7,-1) rectangle +(1,1);
\fill[tintA] (-6,-1) rectangle +(1,1);
\fill[tintA] (-5,-1) rectangle +(1,1);
\fill[tintO] (-4,-1) rectangle +(1,1);
\fill[tintO] (-3,-1) rectangle +(1,1);
\fill[tintO] (-2,-1) rectangle +(1,1);
\fill[tintO] (-1,-1) rectangle +(1,1);
\fill[tintO] (0,-1) rectangle +(1,1);
\fill[tintO] (1,-1) rectangle +(1,1);
\fill[tintO] (2,-1) rectangle +(1,1);
\fill[tintO] (3,-1) rectangle +(1,1);
\fill[tintO] (4,-1) rectangle +(1,1);
\fill[tintO] (5,-1) rectangle +(1,1);
\fill[tintO] (6,-1) rectangle +(1,1);
\fill[tintA] (7,-1) rectangle +(1,1);
\fill[tintA] (8,-1) rectangle +(1,1);
\fill[tintB] (9,-1) rectangle +(1,1);
\fill[tintA] (-9,-2) rectangle +(1,1);
\fill[tintB] (-8,-2) rectangle +(1,1);
\fill[tintB] (-7,-2) rectangle +(1,1);
\fill[tintA] (-6,-2) rectangle +(1,1);
\fill[tintA] (-5,-2) rectangle +(1,1);
\fill[tintO] (-4,-2) rectangle +(1,1);
\fill[tintO] (-3,-2) rectangle +(1,1);
\fill[tintO] (-2,-2) rectangle +(1,1);
\fill[tintO] (-1,-2) rectangle +(1,1);
\fill[tintO] (0,-2) rectangle +(1,1);
\fill[tintO] (1,-2) rectangle +(1,1);
\fill[tintO] (2,-2) rectangle +(1,1);
\fill[tintO] (3,-2) rectangle +(1,1);
\fill[tintO] (4,-2) rectangle +(1,1);
\fill[tintO] (5,-2) rectangle +(1,1);
\fill[tintA] (6,-2) rectangle +(1,1);
\fill[tintA] (7,-2) rectangle +(1,1);
\fill[tintB] (8,-2) rectangle +(1,1);
\fill[tintB] (9,-2) rectangle +(1,1);
\fill[tintA] (-9,-3) rectangle +(1,1);
\fill[tintA] (-8,-3) rectangle +(1,1);
\fill[tintB] (-7,-3) rectangle +(1,1);
\fill[tintB] (-6,-3) rectangle +(1,1);
\fill[tintA] (-5,-3) rectangle +(1,1);
\fill[tintA] (-4,-3) rectangle +(1,1);
\fill[tintO] (-3,-3) rectangle +(1,1);
\fill[tintO] (-2,-3) rectangle +(1,1);
\fill[tintO] (-1,-3) rectangle +(1,1);
\fill[tintO] (0,-3) rectangle +(1,1);
\fill[tintO] (1,-3) rectangle +(1,1);
\fill[tintO] (2,-3) rectangle +(1,1);
\fill[tintO] (3,-3) rectangle +(1,1);
\fill[tintO] (4,-3) rectangle +(1,1);
\fill[tintA] (5,-3) rectangle +(1,1);
\fill[tintA] (6,-3) rectangle +(1,1);
\fill[tintB] (7,-3) rectangle +(1,1);
\fill[tintB] (8,-3) rectangle +(1,1);
\fill[tintA] (9,-3) rectangle +(1,1);
\fill[tintB] (-9,-4) rectangle +(1,1);
\fill[tintA] (-8,-4) rectangle +(1,1);
\fill[tintA] (-7,-4) rectangle +(1,1);
\fill[tintB] (-6,-4) rectangle +(1,1);
\fill[tintB] (-5,-4) rectangle +(1,1);
\fill[tintA] (-4,-4) rectangle +(1,1);
\fill[tintA] (-3,-4) rectangle +(1,1);
\fill[tintO] (-2,-4) rectangle +(1,1);
\fill[tintO] (-1,-4) rectangle +(1,1);
\fill[tintO] (0,-4) rectangle +(1,1);
\fill[tintO] (1,-4) rectangle +(1,1);
\fill[tintO] (2,-4) rectangle +(1,1);
\fill[tintO] (3,-4) rectangle +(1,1);
\fill[tintA] (4,-4) rectangle +(1,1);
\fill[tintA] (5,-4) rectangle +(1,1);
\fill[tintB] (6,-4) rectangle +(1,1);
\fill[tintB] (7,-4) rectangle +(1,1);
\fill[tintA] (8,-4) rectangle +(1,1);
\fill[tintA] (9,-4) rectangle +(1,1);
\fill[tintB] (-9,-5) rectangle +(1,1);
\fill[tintB] (-8,-5) rectangle +(1,1);
\fill[tintA] (-7,-5) rectangle +(1,1);
\fill[tintA] (-6,-5) rectangle +(1,1);
\fill[tintB] (-5,-5) rectangle +(1,1);
\fill[tintB] (-4,-5) rectangle +(1,1);
\fill[tintA] (-3,-5) rectangle +(1,1);
\fill[tintA] (-2,-5) rectangle +(1,1);
\fill[tintO] (-1,-5) rectangle +(1,1);
\fill[tintO] (0,-5) rectangle +(1,1);
\fill[tintO] (1,-5) rectangle +(1,1);
\fill[tintO] (2,-5) rectangle +(1,1);
\fill[tintA] (3,-5) rectangle +(1,1);
\fill[tintA] (4,-5) rectangle +(1,1);
\fill[tintB] (5,-5) rectangle +(1,1);
\fill[tintB] (6,-5) rectangle +(1,1);
\fill[tintA] (7,-5) rectangle +(1,1);
\fill[tintA] (8,-5) rectangle +(1,1);
\fill[tintB] (9,-5) rectangle +(1,1);
\fill[tintA] (-9,-6) rectangle +(1,1);
\fill[tintB] (-8,-6) rectangle +(1,1);
\fill[tintB] (-7,-6) rectangle +(1,1);
\fill[tintA] (-6,-6) rectangle +(1,1);
\fill[tintA] (-5,-6) rectangle +(1,1);
\fill[tintB] (-4,-6) rectangle +(1,1);
\fill[tintB] (-3,-6) rectangle +(1,1);
\fill[tintA] (-2,-6) rectangle +(1,1);
\fill[tintA] (-1,-6) rectangle +(1,1);
\fill[tintA] (0,-6) rectangle +(1,1);
\fill[tintA] (1,-6) rectangle +(1,1);
\fill[tintA] (2,-6) rectangle +(1,1);
\fill[tintA] (3,-6) rectangle +(1,1);
\fill[tintB] (4,-6) rectangle +(1,1);
\fill[tintB] (5,-6) rectangle +(1,1);
\fill[tintA] (6,-6) rectangle +(1,1);
\fill[tintA] (7,-6) rectangle +(1,1);
\fill[tintB] (8,-6) rectangle +(1,1);
\fill[tintB] (9,-6) rectangle +(1,1);
\fill[tintA] (-9,-7) rectangle +(1,1);
\fill[tintA] (-8,-7) rectangle +(1,1);
\fill[tintB] (-7,-7) rectangle +(1,1);
\fill[tintB] (-6,-7) rectangle +(1,1);
\fill[tintA] (-5,-7) rectangle +(1,1);
\fill[tintA] (-4,-7) rectangle +(1,1);
\fill[tintB] (-3,-7) rectangle +(1,1);
\fill[tintB] (-2,-7) rectangle +(1,1);
\fill[tintA] (-1,-7) rectangle +(1,1);
\fill[tintA] (0,-7) rectangle +(1,1);
\fill[tintA] (1,-7) rectangle +(1,1);
\fill[tintA] (2,-7) rectangle +(1,1);
\fill[tintB] (3,-7) rectangle +(1,1);
\fill[tintB] (4,-7) rectangle +(1,1);
\fill[tintA] (5,-7) rectangle +(1,1);
\fill[tintA] (6,-7) rectangle +(1,1);
\fill[tintB] (7,-7) rectangle +(1,1);
\fill[tintB] (8,-7) rectangle +(1,1);
\fill[tintA] (9,-7) rectangle +(1,1);
\fill[tintB] (-9,-8) rectangle +(1,1);
\fill[tintA] (-8,-8) rectangle +(1,1);
\fill[tintA] (-7,-8) rectangle +(1,1);
\fill[tintB] (-6,-8) rectangle +(1,1);
\fill[tintB] (-5,-8) rectangle +(1,1);
\fill[tintA] (-4,-8) rectangle +(1,1);
\fill[tintA] (-3,-8) rectangle +(1,1);
\fill[tintB] (-2,-8) rectangle +(1,1);
\fill[tintB] (-1,-8) rectangle +(1,1);
\fill[tintB] (0,-8) rectangle +(1,1);
\fill[tintB] (1,-8) rectangle +(1,1);
\fill[tintB] (2,-8) rectangle +(1,1);
\fill[tintB] (3,-8) rectangle +(1,1);
\fill[tintA] (4,-8) rectangle +(1,1);
\fill[tintA] (5,-8) rectangle +(1,1);
\fill[tintB] (6,-8) rectangle +(1,1);
\fill[tintB] (7,-8) rectangle +(1,1);
\fill[tintA] (8,-8) rectangle +(1,1);
\fill[tintA] (9,-8) rectangle +(1,1);
\fill[tintB] (-9,-9) rectangle +(1,1);
\fill[tintB] (-8,-9) rectangle +(1,1);
\fill[tintA] (-7,-9) rectangle +(1,1);
\fill[tintA] (-6,-9) rectangle +(1,1);
\fill[tintB] (-5,-9) rectangle +(1,1);
\fill[tintB] (-4,-9) rectangle +(1,1);
\fill[tintA] (-3,-9) rectangle +(1,1);
\fill[tintA] (-2,-9) rectangle +(1,1);
\fill[tintB] (-1,-9) rectangle +(1,1);
\fill[tintB] (0,-9) rectangle +(1,1);
\fill[tintB] (1,-9) rectangle +(1,1);
\fill[tintB] (2,-9) rectangle +(1,1);
\fill[tintA] (3,-9) rectangle +(1,1);
\fill[tintA] (4,-9) rectangle +(1,1);
\fill[tintB] (5,-9) rectangle +(1,1);
\fill[tintB] (6,-9) rectangle +(1,1);
\fill[tintA] (7,-9) rectangle +(1,1);
\fill[tintA] (8,-9) rectangle +(1,1);
\fill[tintB] (9,-9) rectangle +(1,1);
\draw[hair,step=1] (-9,-9) grid (10,10);
\node at (-8.5,9.5) {10};
\node at (-7.5,9.5) {9};
\node at (-6.5,9.5) {9};
\node at (-5.5,9.5) {9};
\node at (-4.5,9.5) {9};
\node at (-3.5,9.5) {9};
\node at (-2.5,9.5) {9};
\node at (-1.5,9.5) {9};
\node at (-0.5,9.5) {9};
\node at (0.5,9.5) {9};
\node at (1.5,9.5) {9};
\node at (2.5,9.5) {9};
\node at (3.5,9.5) {9};
\node at (4.5,9.5) {9};
\node at (5.5,9.5) {9};
\node at (6.5,9.5) {10};
\node at (7.5,9.5) {10};
\node at (8.5,9.5) {11};
\node at (9.5,9.5) {10};
\node at (-8.5,8.5) {9};
\node at (-7.5,8.5) {9};
\node at (-6.5,8.5) {8};
\node at (-5.5,8.5) {8};
\node at (-4.5,8.5) {8};
\node at (-3.5,8.5) {8};
\node at (-2.5,8.5) {8};
\node at (-1.5,8.5) {8};
\node at (-0.5,8.5) {8};
\node at (0.5,8.5) {8};
\node at (1.5,8.5) {8};
\node at (2.5,8.5) {7};
\node at (3.5,8.5) {9};
\node at (4.5,8.5) {8};
\node at (5.5,8.5) {9};
\node at (6.5,8.5) {8};
\node at (7.5,8.5) {9};
\node at (8.5,8.5) {9};
\node at (9.5,8.5) {10};
\node at (-8.5,7.5) {9};
\node at (-7.5,7.5) {8};
\node at (-6.5,7.5) {8};
\node at (-5.5,7.5) {7};
\node at (-4.5,7.5) {7};
\node at (-3.5,7.5) {7};
\node at (-2.5,7.5) {7};
\node at (-1.5,7.5) {7};
\node at (-0.5,7.5) {7};
\node at (0.5,7.5) {7};
\node at (1.5,7.5) {7};
\node at (2.5,7.5) {7};
\node at (3.5,7.5) {7};
\node at (4.5,7.5) {7};
\node at (5.5,7.5) {7};
\node at (6.5,7.5) {8};
\node at (7.5,7.5) {8};
\node at (8.5,7.5) {9};
\node at (9.5,7.5) {8};
\node at (-8.5,6.5) {8};
\node at (-7.5,6.5) {8};
\node at (-6.5,6.5) {7};
\node at (-5.5,6.5) {7};
\node at (-4.5,6.5) {6};
\node at (-3.5,6.5) {6};
\node at (-2.5,6.5) {6};
\node at (-1.5,6.5) {6};
\node at (-0.5,6.5) {6};
\node at (0.5,6.5) {6};
\node at (1.5,6.5) {6};
\node at (2.5,6.5) {5};
\node at (3.5,6.5) {7};
\node at (4.5,6.5) {6};
\node at (5.5,6.5) {7};
\node at (6.5,6.5) {6};
\node at (7.5,6.5) {7};
\node at (8.5,6.5) {7};
\node at (9.5,6.5) {8};
\node at (-8.5,5.5) {8};
\node at (-7.5,5.5) {7};
\node at (-6.5,5.5) {7};
\node at (-5.5,5.5) {6};
\node at (-4.5,5.5) {6};
\node at (-3.5,5.5) {5};
\node at (-2.5,5.5) {5};
\node at (-1.5,5.5) {5};
\node at (-0.5,5.5) {5};
\node at (0.5,5.5) {5};
\node at (1.5,5.5) {5};
\node at (2.5,5.5) {5};
\node at (3.5,5.5) {5};
\node at (4.5,5.5) {5};
\node at (5.5,5.5) {5};
\node at (6.5,5.5) {6};
\node at (7.5,5.5) {6};
\node at (8.5,5.5) {7};
\node at (9.5,5.5) {6};
\node at (-8.5,4.5) {7};
\node at (-7.5,4.5) {7};
\node at (-6.5,4.5) {6};
\node at (-5.5,4.5) {6};
\node at (-4.5,4.5) {5};
\node at (-3.5,4.5) {5};
\node at (-2.5,4.5) {4};
\node at (-1.5,4.5) {4};
\node at (-0.5,4.5) {4};
\node at (0.5,4.5) {4};
\node at (1.5,4.5) {4};
\node at (3.5,4.5) {5};
\node at (4.5,4.5) {4};
\node at (5.5,4.5) {5};
\node at (6.5,4.5) {4};
\node at (7.5,4.5) {5};
\node at (8.5,4.5) {5};
\node at (9.5,4.5) {6};
\node at (-8.5,3.5) {7};
\node at (-7.5,3.5) {6};
\node at (-6.5,3.5) {6};
\node at (-5.5,3.5) {5};
\node at (-4.5,3.5) {5};
\node at (-3.5,3.5) {4};
\node at (-2.5,3.5) {4};
\node at (6.5,3.5) {4};
\node at (7.5,3.5) {4};
\node at (8.5,3.5) {5};
\node at (9.5,3.5) {5};
\node at (-8.5,2.5) {6};
\node at (-7.5,2.5) {6};
\node at (-6.5,2.5) {5};
\node at (-5.5,2.5) {5};
\node at (-4.5,2.5) {4};
\node at (-3.5,2.5) {4};
\node at (7.5,2.5) {4};
\node at (8.5,2.5) {4};
\node at (9.5,2.5) {5};
\node at (-8.5,1.5) {6};
\node at (-7.5,1.5) {5};
\node at (-6.5,1.5) {5};
\node at (-5.5,1.5) {4};
\node at (-4.5,1.5) {4};
\node at (7.5,1.5) {4};
\node at (8.5,1.5) {4};
\node at (9.5,1.5) {5};
\node at (-8.5,0.5) {6};
\node at (-7.5,0.5) {5};
\node at (-6.5,0.5) {5};
\node at (-5.5,0.5) {4};
\node at (-4.5,0.5) {4};
\node at (7.5,0.5) {4};
\node at (8.5,0.5) {4};
\node at (9.5,0.5) {5};
\node at (-8.5,-0.5) {6};
\node at (-7.5,-0.5) {5};
\node at (-6.5,-0.5) {5};
\node at (-5.5,-0.5) {4};
\node at (-4.5,-0.5) {4};
\node at (7.5,-0.5) {4};
\node at (8.5,-0.5) {4};
\node at (9.5,-0.5) {5};
\node at (-8.5,-1.5) {6};
\node at (-7.5,-1.5) {5};
\node at (-6.5,-1.5) {5};
\node at (-5.5,-1.5) {4};
\node at (-4.5,-1.5) {4};
\node at (6.5,-1.5) {4};
\node at (7.5,-1.5) {4};
\node at (8.5,-1.5) {5};
\node at (9.5,-1.5) {5};
\node at (-8.5,-2.5) {6};
\node at (-7.5,-2.5) {6};
\node at (-6.5,-2.5) {5};
\node at (-5.5,-2.5) {5};
\node at (-4.5,-2.5) {4};
\node at (-3.5,-2.5) {4};
\node at (5.5,-2.5) {4};
\node at (6.5,-2.5) {4};
\node at (7.5,-2.5) {5};
\node at (8.5,-2.5) {5};
\node at (9.5,-2.5) {6};
\node at (-8.5,-3.5) {7};
\node at (-7.5,-3.5) {6};
\node at (-6.5,-3.5) {6};
\node at (-5.5,-3.5) {5};
\node at (-4.5,-3.5) {5};
\node at (-3.5,-3.5) {4};
\node at (-2.5,-3.5) {4};
\node at (4.5,-3.5) {4};
\node at (5.5,-3.5) {4};
\node at (6.5,-3.5) {5};
\node at (7.5,-3.5) {5};
\node at (8.5,-3.5) {6};
\node at (9.5,-3.5) {6};
\node at (-8.5,-4.5) {7};
\node at (-7.5,-4.5) {7};
\node at (-6.5,-4.5) {6};
\node at (-5.5,-4.5) {6};
\node at (-4.5,-4.5) {5};
\node at (-3.5,-4.5) {5};
\node at (-2.5,-4.5) {4};
\node at (-1.5,-4.5) {4};
\node at (3.5,-4.5) {4};
\node at (4.5,-4.5) {4};
\node at (5.5,-4.5) {5};
\node at (6.5,-4.5) {5};
\node at (7.5,-4.5) {6};
\node at (8.5,-4.5) {6};
\node at (9.5,-4.5) {7};
\node at (-8.5,-5.5) {8};
\node at (-7.5,-5.5) {7};
\node at (-6.5,-5.5) {7};
\node at (-5.5,-5.5) {6};
\node at (-4.5,-5.5) {6};
\node at (-3.5,-5.5) {5};
\node at (-2.5,-5.5) {5};
\node at (-1.5,-5.5) {4};
\node at (-0.5,-5.5) {4};
\node at (0.5,-5.5) {4};
\node at (1.5,-5.5) {4};
\node at (2.5,-5.5) {4};
\node at (3.5,-5.5) {4};
\node at (4.5,-5.5) {5};
\node at (5.5,-5.5) {5};
\node at (6.5,-5.5) {6};
\node at (7.5,-5.5) {6};
\node at (8.5,-5.5) {7};
\node at (9.5,-5.5) {7};
\node at (-8.5,-6.5) {8};
\node at (-7.5,-6.5) {8};
\node at (-6.5,-6.5) {7};
\node at (-5.5,-6.5) {7};
\node at (-4.5,-6.5) {6};
\node at (-3.5,-6.5) {6};
\node at (-2.5,-6.5) {5};
\node at (-1.5,-6.5) {5};
\node at (-0.5,-6.5) {4};
\node at (0.5,-6.5) {4};
\node at (1.5,-6.5) {4};
\node at (2.5,-6.5) {4};
\node at (3.5,-6.5) {5};
\node at (4.5,-6.5) {5};
\node at (5.5,-6.5) {6};
\node at (6.5,-6.5) {6};
\node at (7.5,-6.5) {7};
\node at (8.5,-6.5) {7};
\node at (9.5,-6.5) {8};
\node at (-8.5,-7.5) {9};
\node at (-7.5,-7.5) {8};
\node at (-6.5,-7.5) {8};
\node at (-5.5,-7.5) {7};
\node at (-4.5,-7.5) {7};
\node at (-3.5,-7.5) {6};
\node at (-2.5,-7.5) {6};
\node at (-1.5,-7.5) {5};
\node at (-0.5,-7.5) {5};
\node at (0.5,-7.5) {5};
\node at (1.5,-7.5) {5};
\node at (2.5,-7.5) {5};
\node at (3.5,-7.5) {5};
\node at (4.5,-7.5) {6};
\node at (5.5,-7.5) {6};
\node at (6.5,-7.5) {7};
\node at (7.5,-7.5) {7};
\node at (8.5,-7.5) {8};
\node at (9.5,-7.5) {8};
\node at (-8.5,-8.5) {9};
\node at (-7.5,-8.5) {9};
\node at (-6.5,-8.5) {8};
\node at (-5.5,-8.5) {8};
\node at (-4.5,-8.5) {7};
\node at (-3.5,-8.5) {7};
\node at (-2.5,-8.5) {6};
\node at (-1.5,-8.5) {6};
\node at (-0.5,-8.5) {5};
\node at (0.5,-8.5) {5};
\node at (1.5,-8.5) {5};
\node at (2.5,-8.5) {5};
\node at (3.5,-8.5) {6};
\node at (4.5,-8.5) {6};
\node at (5.5,-8.5) {7};
\node at (6.5,-8.5) {7};
\node at (7.5,-8.5) {8};
\node at (8.5,-8.5) {8};
\node at (9.5,-8.5) {9};
\node[anchor=north] at (-8.5,-9.35) {$-9$};
\node[anchor=east] at (-9.35,-8.5) {$-9$};
\node[anchor=north] at (-5.5,-9.35) {$-6$};
\node[anchor=east] at (-9.35,-5.5) {$-6$};
\node[anchor=north] at (-2.5,-9.35) {$-3$};
\node[anchor=east] at (-9.35,-2.5) {$-3$};
\node[anchor=north] at (0.5,-9.35) {$0$};
\node[anchor=east] at (-9.35,0.5) {$0$};
\node[anchor=north] at (3.5,-9.35) {$3$};
\node[anchor=east] at (-9.35,3.5) {$3$};
\node[anchor=north] at (6.5,-9.35) {$6$};
\node[anchor=east] at (-9.35,6.5) {$6$};
\node[anchor=north] at (9.5,-9.35) {$9$};
\node[anchor=east] at (-9.35,9.5) {$9$};
\draw[accent,line width=1.3pt] (0.06,0.06) rectangle +(0.88,0.88);
\end{tikzpicture}
\caption{The lap number \(\lambda(u,v)\) for \(|u|, |v| \le 9\), with \(u\) to the right and \(v\) upwards.
Even and odd laps from 4 on have different tints. The 80 squares with \(\lambda \le 3\) are grey and
carry no number, and the origin is outlined.}
\end{figure}

\begin{lemma}[the first 80 squares]
\unbroken
The squares \(T_0, \dots, T_{79}\) are distinct. A square \(s\) has \(\lambda(s) \le 3\) exactly if
it is one of them.
\end{lemma}

\begin{lemma}[indices]
\unbroken
Let \(s\) be a square with \(\lambda(s) = R \ge 4\). Then \(\mathrm{lo}(R) \le \kappa(s) <
\mathrm{hi}(R)\) and \(P_R(\kappa(s)) = s\).
\end{lemma}

\begin{example}
\unbroken
The square \(s = (-5,-2)\) lies in region 4. Tables 2 and 5 give
\[ \lambda(s) = \lfloor (1+5+2)/2 \rfloor = 4, \qquad \kappa(s) = 4 \cdot 4 - 2 - 5 + 2 = 11. \]
In Table~1 with \(R = 4\), the index \(11 = 2R + 2 + 1\) belongs to piece D with \(t = 1\), and
\[ P_4(11) = (2 - 8 + \lfloor 2/2 \rfloor,\ -2 - \lfloor 1/2 \rfloor) = (-5,-2). \]
\end{example}

\begin{definition}[stamps]
\unbroken
The \emph{stamp} of a square \(s\) is the pair
\[ \tau(s) = \begin{cases} (3,\ i) & \text{if } s = T_i, \\ (\lambda(s),\ \kappa(s)) & \text{if } \lambda(s) \ge 4. \end{cases} \]
Stamps are ordered lexicographically. Put \(\mathrm{lo}(3) = 0\) and \(\mathrm{hi}(3) = 80\), and
let \(S\) be the set of pairs \((R,j)\) of integers with \(R \ge 3\) and \(\mathrm{lo}(R) \le j <
\mathrm{hi}(R)\). The \emph{successor} of \((R,j)\) in \(S\) is \((R, j+1)\) if \(j + 1 <
\mathrm{hi}(R)\), and \((R+1, \mathrm{lo}(R+1))\) otherwise.
\end{definition}

\noindent
By Lemma~2 every square has a stamp.

\begin{corollary}
\unbroken
Every stamp lies in \(S\), and different squares have different stamps.
\end{corollary}

\begin{proof}
The first claim is Lemma~2 for stamps \((3,i)\) and Lemma~3 for the others. Let \(\tau(s) =
\tau(s')\). If this stamp is \((3,i)\), then \(s = T_i = s'\). If it is \((R,j)\) with \(R \ge 4\),
then \(s = P_R(j) = s'\) by Lemma~3.
\end{proof}

\begin{definition}[next square]
\unbroken
The \emph{next square} \(s^+\) of a square \(s\) is
\[ s^+ = \begin{cases}
 T_{i+1} & \text{if } s = T_i \text{ with } i < 79, \\
 P_4(\mathrm{lo}(4)) & \text{if } s = T_{79}, \\
 P_R(\kappa(s)+1) & \text{if } \lambda(s) = R \ge 4 \text{ and } \kappa(s) + 1 < \mathrm{hi}(R), \\
 P_{R+1}(\mathrm{lo}(R+1)) & \text{if } \lambda(s) = R \ge 4 \text{ and } \kappa(s) + 1 = \mathrm{hi}(R).
\end{cases} \]
\end{definition}

\noindent
By Lemmas 2 and 3 exactly one case applies.

\begin{lemma}[the next square]
\unbroken
For every square \(s\), the difference \(s^+ - s\) is a step, and \(\tau(s^+)\) is the successor of
\(\tau(s)\) in \(S\).
\end{lemma}

\begin{lemma}[the greedy rule]
\unbroken
Let \(s\) be a square and \(e\) a step. If \(\tau(s+e) > \tau(s)\), then \(\sigma(s^+) \le
\sigma(s+e)\).
\end{lemma}

\noindent
In words: a square one step from \(s\) with a later stamp never has a lower number than the next
square \(s^+\).

\emph{On the proofs of Lemmas 2, 3, 5 and 6.} We do not print them. A case check on Table~2 shows
that all squares with \(\lambda \le 3\) lie in the box \(-4 \le u \le 6\), \(-5 \le v \le 4\). For
them the lemmas are checked by evaluation. For a square with \(\lambda(s) \ge 4\) the proof splits
into cases: by the region of \(s\), by the rows of Tables 2 and 5, by the remainders these rows use,
and in Lemmas 5 and 6 also by the region of the second square. In every case each quantity is an
explicit expression in \(u\) and \(v\) that is linear apart from division by 2, 3 or 6 with
remainder. So every case is a statement of linear arithmetic over the integers, and such statements
can be decided mechanically. Lemma~6 compares spiral numbers, which are quadratic. By Lemma~1~(iii)
it is enough to compare rings and positions, and these are linear. In Lean every case is closed by
\texttt{omega}, a decision procedure for linear integer arithmetic, and the statements about the
first 80 squares by \texttt{decide}, which evaluates them inside the kernel.

\section{Proof of Theorems 1 and 2}

Let \(V_0 = (0,0) = T_0\) and \(V_{n+1} = V_n^+\).

\emph{Step 1. The stamps \(\tau(V_0), \tau(V_1), \dots\) are the elements of \(S\) in increasing
order, each once.} For every \(R \ge 3\) the number \(\mathrm{hi}(R) - \mathrm{lo}(R)\) is positive.
So \((3,0)\) is the least element of \(S\), and every element of \(S\) is reached from \((3,0)\) by
taking the successor finitely often. Now \(\tau(V_0) = (3,0)\), and \(\tau(V_{n+1})\) is the
successor of \(\tau(V_n)\) by Lemma~5.

\emph{Step 2. \(V_n = W_n\) for all \(n\), and Theorem 2 (a) holds.} Let \(s\) be a square. By
Corollary~4 and Step~1 there is exactly one \(n\) with \(\tau(V_n) = \tau(s)\), and then \(V_n = s\)
because different squares have different stamps. So every square occurs exactly once among the
\(V_n\). The square with stamp \((3,i)\) is \(T_i\). The square with stamp \((R,j)\), \(R \ge 4\),
is \(P_R(j)\) by Lemma~3. Since the stamps come in increasing order, the sequence \(V\) is \(T_0,
\dots, T_{79}\), then lap 4, then lap 5, and so on. This is the sequence \(W\).

\emph{Step 3. \(W\) satisfies (W1) to (W4), with ``differs by a step'' for ``one move apart''.} (W1)
holds by definition, (W2) by Step~2 and (W3) by Lemma~5. For (W4) let \(e\) be a step and \(q = W_n
+ e\) with \(\sigma(q) < \sigma(W_{n+1})\). Since \(W_{n+1} = W_n^+\), Lemma~6 shows that \(\tau(q)
> \tau(W_n)\) is false. So \(\tau(q) \le \tau(W_n)\). By Step~2, \(q = W_m\) for some \(m\), and by
Step~1, \(m \le n\).

By Section~2 the sequence \(a_n = (u_n + v_n,\ u_n - v_n)\) is therefore a Falafel walk. This proves
Theorem~1~(a). By Remark~1 it is the only Falafel walk, which is Theorem~2~(b). By Step~2 it visits
every square with \(x+y\) even, and by Section~2 it visits no other square. This is
Theorem~1~(b).\qed

\section{The sequence}

In Lemma~7 and Corollary~8 a square is again a pair \((x,y)\), with any parity of \(x+y\).

\begin{lemma}
\unbroken
For every square \((x,y)\), \(\sigma(x,y) \equiv x + y + 1 \pmod 2\).
\end{lemma}

\begin{proof}
Each square of the spiral shares a side with the next one. So from one square to the next the number
goes up by 1 and \(x + y\) changes by 1, and \(\sigma(x,y) - x - y\) keeps its parity. At the square
\((0,0)\) it is 1.
\end{proof}

\begin{corollary}
\unbroken
Let \(a_0, a_1, a_2, \dots\) be the Falafel walk. Every odd positive integer is \(\sigma(a_n)\) for
exactly one \(n\), and every \(\sigma(a_n)\) is odd. So A399160 is a permutation of the odd positive
integers.
\end{corollary}

\begin{proof}
Let \(k\) be an odd positive integer. It is the number of a square \((x,y)\), and \(x + y\) is even
by Lemma~7. By Theorem~1~(b) there is \(n\) with \(a_n = (x,y)\), so \(\sigma(a_n) = k\). If also
\(\sigma(a_m) = k\), then \(a_m = a_n\), because different squares have different numbers, and \(m =
n\) by (W2). Finally, every \(a_n\) has \(x + y\) even, so every \(\sigma(a_n)\) is odd by Lemma~7.
\end{proof}

\begin{corollary}[the terms]
\unbroken
For an integer \(R \ge 4\) put
\[ \mathrm{start}(R) = \begin{cases} (27R^2 - 22R - 24)/4 & \text{if } R \text{ even} \\ (27R^2 - 22R - 33)/4 & \text{if } R \text{ odd} \end{cases} \]
Lap \(R\) starts with \(W_{\mathrm{start}(R)}\). So \(W_n = T_n\) for \(n < 80\). For \(n \ge 80\)
let \(R\) be the largest integer with \(\mathrm{start}(R) \le n\). Then \(W_n\) lies in lap \(R\),
and
\[ W_n = P_R\bigl(\mathrm{lo}(R) + n - \mathrm{start}(R)\bigr). \]
The \((n+1)\)-th term of A399160 is \((2\rho(W_n) - 1)^2 + \pi(W_n)\). Table~6 in Appendix~A gives
this number as a function of \(R\) and the index \(j = \mathrm{lo}(R) + n - \mathrm{start}(R)\).
\end{corollary}

\begin{proof}
By Theorem~2 the squares of the walk are \(W_0, W_1, W_2, \dots\), and Lemma~1~(i) gives the number
of \(W_n\). Lap 4 starts with \(W_{80}\), and \(\mathrm{start}(4) = 80\). For both parities of
\(R\), the difference \(\mathrm{start}(R+1) - \mathrm{start}(R)\) is the number \(\mathrm{hi}(R) -
\mathrm{lo}(R)\) of squares of lap \(R\) (Section~3). So lap \(R\) starts with
\(W_{\mathrm{start}(R)}\) and ends just before \(W_{\mathrm{start}(R+1)}\). The numbers
\(\mathrm{start}(R)\) increase with \(R\). So \(W_n\) with \(n \ge 80\) lies in lap \(R\) exactly if
\(R\) is the largest integer with \(\mathrm{start}(R) \le n\), and then it is the square of lap
\(R\) at index \(\mathrm{lo}(R) + n - \mathrm{start}(R)\). On every row of Table~1 the ring \(\rho\)
and the position \(\pi\) are linear in \(R\) and \(j\), and Table~6 is \((2\rho - 1)^2 + \pi\) for
that row.
\end{proof}

The first 80 terms, the numbers of \(T_0, \dots, T_{79}\), are 1, 3, 7, 19, 5, 9, 11, 15, 33, 13,
29, 53, 25, 27, 47, 23, 45, 21, 41, 69, 17, 35, 39, 67, 37, 63, 97, 61, 95, 31, 55, 59, 93, 57, 89,
129, 87, 127, 51, 49, 79, 117, 77, 115, 43, 71, 75, 113, 73, 109, 153, 107, 151, 105, 149, 65, 99,
103, 147, 101, 143, 193, 141, 191, 139, 189, 137, 187, 91, 131, 135, 185, 133, 181, 237, 179, 235,
177, 85, 125.

\section{The formal proof}

The proof has been formalised in Lean~4, version 4.36.0, with Mathlib. The Lean kernel accepts it.
It has no \texttt{sorry}, which would mark an unproved step, and no \texttt{native\_decide}, which
would take the result of compiled code on trust. It depends only on the three standard axioms
\texttt{propext}, \texttt{Classical.choice} and \texttt{Quot.sound}. The Lean files, with the
program that printed the tables of this paper, are in the repository~\cite{ref5}.

\emph{What a reader has to trust.} Appendix~B reproduces the statement as Lean has it. It defines
\(\sigma\) by formula (1), a move, and a Falafel walk as in Definition~1, and it states four claims:
Remark~1, Theorem~1~(a), Theorem~1~(b) and Corollary~8. The laps, the lap number, the index and the
stamps are only tools of the proof. A mistake in them could make the proof fail. It could not make a
false theorem pass. Corollary~9 is proved in Lean as a formula for the terms: for every \(n \ge 80\)
the number of \(a_n\) is the entry of Table~6 (\texttt{term\_formula}). Its statement is in a second
file, which defines \(\mathrm{start}\), \(\mathrm{lo}\) and the numbers of Table~6 and mentions no
square and no lap. Appendix~B shows it without the rows of the table. A reader who wants that claim
has to trust that file as well.

\emph{What corresponds to what.} Table~3 lists the Lean theorem behind each statement of this paper.
Lean also proves that the first 70 squares of its walk have the 70 numbers listed in the OEIS entry
(\texttt{first\_terms}). This checks that Definition~1 is the rule of the OEIS entry.

\begin{table}[htbp]
\caption{The statements of this paper and the Lean theorems behind them.}
\centering\small
\begin{tabular}{@{}l>{\raggedright\arraybackslash}p{0.72\textwidth}@{}}
\toprule
\textit{Statement} & \textit{Lean theorems} \\
\midrule
The numbering & \texttt{spiral\_uv\_inj} (for squares with x + y even), \texttt{spiral\_pos}, \texttt{spiral\_surjective} \\
Remark 1 & \texttt{deterministic} \\
Theorem 1 (a) & \texttt{never\_trapped} \\
Theorem 1 (b) & \texttt{visits\_exactly\_even\_squares} \\
Theorem 2 & \texttt{first\_eighty}, \texttt{walk\_follows\_path}, \texttt{w\_injective}, \texttt{w\_surjective}, \texttt{walk\_isWalk} \\
Lemma 1 & \texttt{spiral\_uv}, \texttt{rpos\_bounds}, \texttt{spiral\_le\_of\_keyLE} \\
Lemma 2 & \texttt{prefix\_table}, \texttt{prefix\_mem} \\
Lemma 3 & \texttt{pos\_range}, \texttt{sq\_pos} \\
Corollary 4 & \texttt{idx\_range}, \texttt{stamp\_inj} \\
Lemma 5 & \texttt{step\_ok}, \texttt{rook\_ok}, with \texttt{sq\_pos} \\
Lemma 6 & \texttt{greedy} \\
Lemma 7 & \texttt{spiral\_parity} \\
Corollary 8 & \texttt{permutation\_of\_odd\_numbers} \\
Corollary 9 & \texttt{term\_formula}, from \texttt{start\_succ}, \texttt{walk\_eq\_path}, \texttt{ring\_rpos\_core} and \texttt{term\_eq}; for \(n < 80\), \texttt{first\_eighty} \\
\bottomrule
\end{tabular}
\end{table}

\emph{Where Lean and this paper differ.} Lean defines the next square by a table of moves, one move
for every square, arranged by region like Table~5. It proves that this move is a step and leads to
the square with the next stamp. With Lemma~3 this is Lemma~5. Lean also proves that a lap is exactly
a level set of the lap number: for \(R \ge 4\), \(\lambda(s) = R\) exactly if \(s = P_R(j)\) for
some \(j\) with \(\mathrm{lo}(R) \le j < \mathrm{hi}(R)\) (\texttt{lapnum\_eq\_iff}). Lean has the
numbering only as formula (1). So it also proves, as far as the proof needs it, what this paper
takes from the construction of the spiral: different squares with \(x+y\) even have different
numbers, every number is positive, and every positive integer is the number of a square (the first
row of Table~3). It proves Lemma~7 from the four cases of (1).

\emph{What is not formalised.} The tables of this paper are a transcription of the Lean definitions.
A program printed them from expressions that it first compared with the Lean definitions: Tables 2
and 5, \(\rho\) and \(\pi\) for all squares with \(|u|, |v| \le 300\), and Table~1 with
\(\mathrm{lo}\) and \(\mathrm{hi}\) for the laps \(4 \le R \le 300\). The same program compared the
sequence \(W\) of Theorem~2 with a direct simulation of the rule for the first 200,000 steps, and it
tested Lemmas 2, 3, 5 and 6 in the wording of this paper for all squares with \(|u|, |v| \le 120\).
It also checked that formula (1) gives the numbers of the first 160,801 squares along the spiral.
The description of a lap and of the visited rings in Section~3 is not formalised either. The program
compared Table~6 with the Lean table of the terms and with the numbers of the squares of Table~1 for
\(4 \le R \le 300\), and with the simulation. It compared the formula for \(\mathrm{start}(R)\) with
the Lean definition, which is one expression for both parities of \(R\), and with the lengths of the
laps for \(4 \le R \le 5000\).

\emph{Acknowledgement.} The lap structure was found by experiment. The Lean development was written
with the help of Claude, an AI system made by Anthropic. The correctness of the result rests on the
Lean kernel and on the statement in Appendix~B, not on how the proof was written.

\appendix

\section{The first 80 squares, the index and the terms}

\begin{table}[H]
\setlength{\abovecaptionskip}{0pt}
\caption{The first 80 squares \(T_0, \dots, T_{79}\) of the walk, in the coordinates \((u,v)\). Each row
holds eight squares.}
\vspace{\belowcaptionskip}
\centering\small
\setlength{\tabcolsep}{4.5pt}
\begin{tabular}{@{}r@{\hspace{1.4em}}llllllll@{}}
\toprule
\(i\) & \multicolumn{8}{l}{\(T_i,\ T_{i+1},\ \dots,\ T_{i+7}\)} \\
\midrule
0 & \((0,\,0)\) & \((1,\,0)\) & \((-1,\,0)\) & \((-1,\,-1)\) & \((0,\,-1)\) & \((0,\,1)\) & \((1,\,1)\) & \((1,\,-1)\) \\
8 & \((2,\,-1)\) & \((2,\,0)\) & \((2,\,1)\) & \((2,\,2)\) & \((0,\,2)\) & \((1,\,2)\) & \((-1,\,2)\) & \((-1,\,1)\) \\
16 & \((-2,\,1)\) & \((-2,\,0)\) & \((-2,\,-1)\) & \((-2,\,-2)\) & \((0,\,-2)\) & \((1,\,-2)\) & \((-1,\,-2)\) & \((-1,\,-3)\) \\
24 & \((0,\,-3)\) & \((1,\,-3)\) & \((2,\,-3)\) & \((2,\,-2)\) & \((3,\,-2)\) & \((3,\,0)\) & \((3,\,1)\) & \((3,\,-1)\) \\
32 & \((4,\,-1)\) & \((4,\,0)\) & \((4,\,1)\) & \((4,\,2)\) & \((3,\,2)\) & \((3,\,3)\) & \((1,\,3)\) & \((0,\,3)\) \\
40 & \((-1,\,3)\) & \((-2,\,3)\) & \((-2,\,2)\) & \((-3,\,2)\) & \((-3,\,0)\) & \((-3,\,-1)\) & \((-3,\,1)\) & \((-4,\,1)\) \\
48 & \((-4,\,0)\) & \((-4,\,-1)\) & \((-4,\,-2)\) & \((-3,\,-2)\) & \((-3,\,-3)\) & \((-2,\,-3)\) & \((-2,\,-4)\) & \((0,\,-4)\) \\
56 & \((1,\,-4)\) & \((-1,\,-4)\) & \((-1,\,-5)\) & \((0,\,-5)\) & \((1,\,-5)\) & \((2,\,-5)\) & \((2,\,-4)\) & \((3,\,-4)\) \\
64 & \((3,\,-3)\) & \((4,\,-3)\) & \((4,\,-2)\) & \((5,\,-2)\) & \((5,\,0)\) & \((5,\,1)\) & \((5,\,-1)\) & \((6,\,-1)\) \\
72 & \((6,\,0)\) & \((6,\,1)\) & \((6,\,2)\) & \((5,\,2)\) & \((5,\,3)\) & \((4,\,3)\) & \((2,\,3)\) & \((2,\,4)\) \\
\bottomrule
\end{tabular}
\end{table}

\begin{table}[H]
\setlength{\abovecaptionskip}{0pt}
\caption{The index \(\kappa(u,v)\) of a square with lap number \(R = \lambda(u,v) \ge 4\). The regions are
those of Table~2, and \(q\) and \(m\) are as there.}
\vspace{\belowcaptionskip}
\centering\small
\begin{tabular}{@{}lll@{}}
\toprule
\textit{Region} & \textit{Case} & \(\kappa(u, v)\) \\
\midrule
1. North row & \(v \text{ odd} \text{ or } u \le -1\) & \(-u - v\) \\
 & \(v \text{ even},\ u = 0\) & \(-v - 1\) \\
 & \(v \text{ even},\ u = 1\) & \(-v\) \\
\addlinespace
2. North-west staircase &  & \(-u - v\) \\
\addlinespace
3. West corner & \(u \text{ even}\) & \(2R - v\) \\
 & \(u \text{ odd},\ v = 0\) & \(2R - 4\) \\
 & \(u \text{ odd},\ v = -1\) & \(2R - 3\) \\
 & \(u \text{ odd},\ v = 1\) & \(2R - 2\) \\
\addlinespace
4. South-west staircase &  & \(4R - 2 + u - v\) \\
\addlinespace
5. South corner & \(v \text{ even},\ u = 0\) & \(6R - 5\) \\
 & \(v \text{ even},\ u = 1\) & \(6R - 4\) \\
 & \(v \text{ even},\ u = -1\) & \(6R - 3\) \\
 & \(v \text{ odd}\) & \(6R - 1 + u\) \\
\addlinespace
6. South-east staircase &  & \(8R - 2 + u + v\) \\
\addlinespace
7. East corner & \(u \text{ odd},\ v = 0\) & \(10R - 4\) \\
 & \(u \text{ odd},\ v = 1\) & \(10R - 3\) \\
 & \(u \text{ odd},\ v = -1\) & \(10R - 2\) \\
 & \(u \text{ even}\) & \(10R + v\) \\
\addlinespace
8. North-east staircase &  & \(12R - u + v\) \\
\addlinespace
9. Column 3 & \(v \text{ even}\) & \(-v - 4\) \\
 & \(v \text{ odd}\) & \(-v - 2\) \\
\addlinespace
10. Return & \(m \le 1\) & \(11R + 3v - 3q - 1\) \\
 & \(2 \le m \le 3\) & \(11R + 3v - 3q - 2\) \\
 & \(m \ge 4\) & \(11R + 3v - 3q - 3\) \\
\bottomrule
\end{tabular}
\end{table}

\begin{table}[H]
\setlength{\abovecaptionskip}{0pt}
\caption{The number \(\sigma(P_R(j))\) of the square at index \(j\) of lap \(R \ge 4\). The pieces and the
ranges of \(t\) are those of Table~1, and \(k = \lfloor t/3 \rfloor\). Lean proves that these are
the terms of A399160 from the 81st on (Corollary~9 and Section~7).}
\vspace{\belowcaptionskip}
\centering\footnotesize
\begin{tabular}{@{}lll@{}}
\toprule
\textit{Piece} & \textit{Index \(j\)} & \textit{Term} \\
\midrule
A. North row & even \(R\), \(j = -R - 1\) & \((2R - 1)^2 + 8R\) \\
 & even \(R\), \(j = -R\) & \((2R + 1)^2 + 2\) \\
 & odd \(R\), \(j = -R - 3\) & \((2R + 3)^2 + 6\) \\
 & odd \(R\), \(j = -R - 2\) & \((2R + 5)^2 + 6\) \\
 & odd \(R\), \(j = -R - 1\) & \((2R + 1)^2 + 2\) \\
 & all other \(j < 0\) & \((4R + 2j - 1)^2 + 14R + 6j\) \\
\addlinespace
B. North-west staircase & \(j = t\), \(t \text{ even}\) & \((4R - 5)^2 + 14R - t - 14\) \\
 & \(t \text{ odd}\) & \((4R - 3)^2 + 14R - t - 7\) \\
\addlinespace
C. West corner & \(j = 2R - 4 + t\), \(t = 0\) & \((4R - 7)^2 + 12R - 18\) \\
 & \(t = 1\) & \((4R - 5)^2 + 12R - 14\) \\
 & \(t = 2\) & \((4R - 5)^2 + 12R - 10\) \\
 & \(t = 3\) & \((4R - 3)^2 + 12R - 4\) \\
 & \(t = 4\) & \((4R - 5)^2 + 12R - 12\) \\
 & \(t = 5\) & \((4R - 3)^2 + 12R - 8\) \\
\addlinespace
D. South-west staircase & \(j = 2R + 2 + t\), \(t \text{ even}\) & \((4R - 1)^2 + 12R - t - 4\) \\
 & \(t \text{ odd}\) & \((4R - 3)^2 + 12R - t - 9\) \\
\addlinespace
E. South corner & \(j = 6R - 5 + t\), \(t = 0\) & \((4R - 5)^2 + 8R - 8\) \\
 & \(t = 1\) & \((4R - 3)^2 + 8R - 6\) \\
 & \(t = 2\) & \((4R - 3)^2 + 8R - 2\) \\
 & \(t = 3\) & \((4R - 1)^2 + 8R + 2\) \\
 & \(t = 4\) & \((4R - 3)^2 + 8R - 4\) \\
 & \(t = 5\) & \((4R - 1)^2 + 8R - 2\) \\
\addlinespace
F. South-east staircase & \(j = 6R + 1 + t\), \(t \text{ even}\) & \((4R + 1)^2 + 8R - t\) \\
 & \(t \text{ odd}\) & \((4R - 1)^2 + 8R - t - 3\) \\
\addlinespace
G. East corner & \(j = 10R - 4 + t\), \(t = 0\) & \((4R - 3)^2 + 4R - 2\) \\
 & \(t = 1\) & \((4R - 1)^2 + 4R - 2\) \\
 & \(t = 2\) & \((4R - 1)^2 + 4R + 2\) \\
 & \(t = 3\) & \((4R + 1)^2 + 4R + 4\) \\
 & \(t = 4\) & \((4R - 1)^2 + 4R\) \\
 & \(t = 5\) & \((4R + 1)^2 + 4R\) \\
\addlinespace
H. North-east staircase & \(j = 10R + 2 + t\), \(t \text{ even}\) & \((4R + 3)^2 + 4R - t\) \\
 & \(t \text{ odd}\) & \((4R + 1)^2 + 4R - t - 1\) \\
\addlinespace
I. Return & \(j = 11R + 3 + t\) &  \\
 & even \(R\), \(t = 3k\) & \((4R - 4k - 1)^2 + 3R - 6k - 4\) \\
 & even \(R\), \(t = 3k + 1\) & \((4R - 4k - 3)^2 + 3R - 6k - 6\) \\
 & even \(R\), \(t = 3k + 2\) & \((4R - 4k - 1)^2 + 3R - 6k - 6\) \\
 & odd \(R\), \(t = 3k\) & \((4R - 4k - 3)^2 + 3R - 6k - 5\) \\
 & odd \(R\), \(t = 3k + 1\) & \((4R - 4k - 1)^2 + 3R - 6k - 5\) \\
 & odd \(R\), \(t = 3k + 2\) & \((4R - 4k - 3)^2 + 3R - 6k - 7\) \\
\bottomrule
\end{tabular}
\end{table}

\clearpage

\section{The statement in Lean}

These are the definitions and the four claims as the Lean proof states them. Each claim is one
statement of this paper and names it in its comment. A square is a pair \texttt{(x, y)} of integers,
and \texttt{a n} is \(a_n\).

\begin{Verbatim}[fontfamily=leanmono, fontsize=\fontsize{7.2pt}{9.4pt}\selectfont, frame=lines, framesep=6pt]
/-- The number of the square (x, y) on the square spiral, as OEIS has it: 1 at (0, 0),
2 at (1, 0), 3 at (1, 1), then counter-clockwise. Here `r` is the ring max(|x|, |y|),
and ring r ≥ 1 holds the numbers (2r - 1)² + 1 to (2r + 1)². -/
def spiral (p : ℤ × ℤ) : ℤ :=
  let x := p.1
  let y := p.2
  let r := max |x| |y|
  if x = r ∧ y > -r then (2 * r - 1) ^ 2 + (y + r)        -- right side, going up
  else if y = r then (2 * r - 1) ^ 2 + 2 * r + (r - x)    -- top side, going left
  else if x = -r then (2 * r - 1) ^ 2 + 4 * r + (r - y)   -- left side, going down
  else (2 * r - 1) ^ 2 + 6 * r + (x + r)                  -- bottom side, going right

/-- The Falafel (Fers + Alfil) moves one or two squares diagonally. -/
def IsMove (p q : ℤ × ℤ) : Prop :=
  |q.1 - p.1| = |q.2 - p.2| ∧ (|q.1 - p.1| = 1 ∨ |q.1 - p.1| = 2)

/-- `a 0, a 1, a 2, ...` is an infinite run of the walk. It starts on square 1 and visits no
square twice. Each next square is one move away. Every square one move away with a lower number
than the next square was visited before. -/
structure IsWalk (a : ℕ → ℤ × ℤ) : Prop where
  start : a 0 = (0, 0)
  injective : Function.Injective a
  move : ∀ n, IsMove (a n) (a (n + 1))
  lowest : ∀ n q, IsMove (a n) q → spiral q < spiral (a (n + 1)) → ∃ m ≤ n, a m = q

/-- Remark 1. The rule decides every step: there is at most one infinite run. -/
def Deterministic : Prop := ∀ a b : ℕ → ℤ × ℤ, IsWalk a → IsWalk b → a = b

/-- Theorem 1 (a). The walk is never trapped: it has an infinite run.
A399160(n) is `spiral (a (n - 1))`. -/
def NeverTrapped : Prop := ∃ a : ℕ → ℤ × ℤ, IsWalk a

/-- Theorem 1 (b). The walk visits every square (x, y) with x + y even, and no other square. -/
def VisitsExactlyEvenSquares : Prop :=
  ∀ a : ℕ → ℤ × ℤ, IsWalk a → ∀ x y : ℤ, (∃ n, a n = (x, y)) ↔ Even (x + y)

/-- Corollary 8. Every term is an odd positive integer, and every odd positive integer is a term
exactly once. So A399160 is a permutation of the odd positive integers. -/
def PermutationOfOddNumbers : Prop :=
  ∀ a : ℕ → ℤ × ℤ, IsWalk a →
    (∀ n, 0 < spiral (a n) ∧ Odd (spiral (a n))) ∧
    (∀ k : ℤ, 0 < k → Odd k → ∃! n, spiral (a n) = k)
\end{Verbatim}

The second file states Corollary~9. This is an excerpt: the file also defines \texttt{backTerm} and
\texttt{term}, and \texttt{term R j} is the entry of Table~6 for lap \(R\) and index \(j\). Here
\texttt{start R} is \(\mathrm{start}(R)\) and \texttt{first R} is \(\mathrm{lo}(R)\).

\begin{Verbatim}[fontfamily=leanmono, fontsize=\fontsize{7.2pt}{9.4pt}\selectfont, frame=lines, framesep=6pt]
/-- The step at which lap R ≥ 4 starts: (27 R² - 22 R - 24)/4 for even R and
(27 R² - 22 R - 33)/4 for odd R. -/
def start (R : ℤ) : ℤ := (27 * R ^ 2 - 22 * R - 24 - 9 * (R % 2)) / 4

/-- Lap 4 starts at step 80, lap 5 at step 133 and lap 6 at step 204. -/
example : start 4 = 80 ∧ start 5 = 133 ∧ start 6 = 204 := by decide

/-- The index of the first step of lap R: -R - 1 for even R and -R - 3 for odd R. -/
def first (R : ℤ) : ℤ := -(R + 1 + 2 * (R % 2))

/-- Corollary 9. If lap R ≥ 4 has started at step n and lap R + 1 has not, then the number of
the square at step n is `term R j` with j = first R + (n - start R). This number is
A399160(n + 1). -/
def TermFormula : Prop :=
  ∀ a : ℕ → ℤ × ℤ, IsWalk a → ∀ R : ℤ, 4 ≤ R → ∀ n : ℕ, start R ≤ n → n < start (R + 1) →
    spiral (a n) = term R (first R + (n - start R))
\end{Verbatim}

\begin{thebibliography}{9}

\bibitem{ref1}
OEIS Foundation Inc., Entry A399160, \emph{The On-Line Encyclopedia of Integer Sequences},
2026.\newline\url{https://oeis.org/A399160}

\bibitem{ref2}
OEIS Foundation Inc., Entry A316667, \emph{The On-Line Encyclopedia of Integer Sequences},
2018.\newline\url{https://oeis.org/A316667}

\bibitem{ref3}
L. de Moura and S. Ullrich, The Lean 4 theorem prover and programming language, in \emph{Automated
Deduction -- CADE 28}, Lecture Notes in Computer Science 12699, Springer, 2021, pp. 625--635.

\bibitem{ref4}
The mathlib Community, The Lean mathematical library, in \emph{Proceedings of the 9th ACM SIGPLAN
International Conference on Certified Programs and Proofs (CPP 2020)}, ACM, 2020, pp. 367--381.

\bibitem{ref5}
M. Norbäck Olivers, \emph{spiral-walks}, source repository,
2026.\newline\url{https://github.com/norpan/spiral-walks}

\end{thebibliography}

\end{document}
